\documentclass[11pt]{article}

\usepackage{amsmath,amssymb,amsthm}
\usepackage[margin=1.1in]{geometry}
\usepackage{bm}
\usepackage{array}
\usepackage[colorlinks=true,linkcolor=blue,citecolor=blue]{hyperref}

\newtheorem{lemma}{Lemma}
\newtheorem{proposition}{Proposition}
\newtheorem{remark}{Remark}
\newcommand{\R}{\mathbb{R}}
\newcommand{\eps}{\varepsilon}
\newcommand{\Bs}{B_{\mathrm{s}}}
\newcommand{\Bf}{B_{\mathrm{f}}}
\newcommand{\bu}{\mathbf{u}}
\newcommand{\bv}{\mathbf{v}}
\newcommand{\bw}{\mathbf{w}}
\newcommand{\bx}{\mathbf{x}}
\newcommand{\bn}{\mathbf{n}}
\newcommand{\bN}{\mathbf{N}}
\newcommand{\bU}{\mathbf{U}}
\newcommand{\bff}{\mathbf{f}}
\newcommand{\br}{\mathbf{r}}
\newcommand{\bfeta}{\bm{\eta}}
\newcommand{\bnu}{\bm{\nu}}
\newcommand{\jump}[1]{[\![#1]\!]}
\newcommand{\dS}{\,\mathrm{d}S}
\newcommand{\dV}{\,\mathrm{d}V}
\newcommand{\dx}{\,\mathrm{d}x}
\newcommand{\tr}{\operatorname{tr}}
\newcommand{\dev}{\operatorname{dev}}
\newcommand{\sym}{\operatorname{sym}}
\newcommand{\cof}{\operatorname{cof}}
\newcommand{\Div}{\operatorname{Div}}
\renewcommand{\div}{\operatorname{div}}
\newcommand{\norm}[1]{\lVert #1 \rVert}
\newcommand{\ip}[2]{\langle #1, #2\rangle}
\newcommand{\Ja}{J_{\mathrm{naive}}}
\newcommand{\JF}{J_{\mathrm{F}}}
\newcommand{\Jmu}{J_{\mu}}
\newcommand{\Jh}{J_{\mathrm{hyd}}}

\title{Equilibrium states and hydrostatic figures by constrained
       optimisation:\\ feasibility functionals, their derivatives,
       and Sobolev gradients}
\author{AdGIA method notes}
\date{}

\begin{document}
\maketitle
\sloppy

\noindent
This note is the theory of computing equilibrium states of
self-gravitating bodies with fluid regions --- and, in its
generalisation, hydrostatic figures --- as the solution of a
constrained optimisation problem. Its companions are
\texttt{doc/gravitating\_elasticity.md} (notation, the relabelling laws,
and section ``Equilibrium stress in general models'', whose
minimum-deviatoric generator is the inner solver used here),
\texttt{doc/slip\_interface.tex} (the slip machinery, which describes the
dynamics \emph{about} an equilibrium, not the equilibrium itself) and
\texttt{doc/gauge\_penalty\_iteration.tex} (gauge handling and its
discrete semi-convergence). The planning summary is
\texttt{doc/planning/equilibrium\_figures.md}.

The organisation: \S\ref{sec:problem} states the equilibrium problem and
shows that the equilibrium mapping may be taken continuous;
\S\ref{sec:functional} defines the feasibility functionals --- a naive
one that fails, an exact fluid-only one, and a weighted whole-body one
that approaches it --- and characterises their zero sets;
\S\ref{sec:dual} derives their common Stokes (dual) structure and, from
it, the exact recovery of the body's stress field;
\S\ref{sec:derivative} derives the derivative with respect to density
and the Hessian action; \S\ref{sec:energy} interprets the functional as
the squared gradient of the gravitational energy and derives the
alternative advection route; \S\ref{sec:sobolev} turns derivatives into
gradients by Sobolev metrics; \S\ref{sec:stage2} admits the shape as a
control; \S\ref{sec:second} derives second derivatives and
Newton--Krylov steps for density and shape; \S\ref{sec:algorithms}
explains the implemented algorithms (nonlinear CG, Levenberg--Marquardt)
and records what has been measured.

\paragraph{Sign convention.} Throughout, as in the rest of the library,
\begin{equation}
  \nabla^{2}\Phi = 4\pi G \varrho \quad\text{in } \R^{3},
  \qquad \Phi \to 0 \ \text{at infinity},
  \label{eq:poisson-strong}
\end{equation}
so that $\Phi < 0$, the gravitational acceleration is $-\nabla\Phi$, the
body force per unit volume is $-\varrho\nabla\Phi$, and equilibrium of a
Cauchy stress $T$ reads $\Div T - \varrho\nabla\Phi = 0$. In a fluid,
$T = -p\,1$ and the hydrostatic balance is $\nabla p = -\varrho\nabla\Phi$.
(With the opposite convention $\Phi > 0$ every body-force term below
changes sign; nothing else changes.) We write
\begin{equation}
  \bff := \varrho\,\nabla\Phi ,
  \qquad\text{so that equilibrium is}\qquad \Div T = \bff .
  \label{eq:f}
\end{equation}
$\Phi$ is always the \emph{total} self-consistent potential of the
whole body (one global Poisson solve; \S\ref{sec:poisson}); in the
rotating case of \S\ref{sec:rotation} a centrifugal potential is
added.

% ======================================================================
\section{The equilibrium problem and the role of the mapping}
\label{sec:problem}

\subsection{Equations at a natural reference}

Let the body occupy $B \subset \R^{3}$, partitioned into a fluid part
$\Bf$ and a solid part $\Bs$, each a union of connected regions, with
material discontinuities across a set of internal interfaces; $\Sigma$
denotes the fluid--solid interfaces and $\partial B$ the free surface.
At a \emph{natural} reference (particle label $=$ equilibrium position)
a static state is a symmetric stress field $T$ and density $\varrho$
such that
\begin{equation}
  \begin{gathered}
  \Div T = \varrho\nabla\Phi \ \text{in each region}, \qquad
  \jump{T\bn} = 0 \ \text{across interfaces}, \\
  T\bn = 0 \ \text{on } \partial B,
  \qquad \dev T = 0 \ \text{in } \Bf ,
  \end{gathered}
  \label{eq:equilibrium}
\end{equation}
with $\Phi$ solving \eqref{eq:poisson-strong} (so that $\Phi$ and
$\nabla\Phi\cdot\bn$ are continuous across every interface). Here
$\dev T = T - \tfrac13(\tr T)1$. The solid stress is otherwise
unconstrained: whatever deviatoric stress the solid carries is, at this
level, admissible (consistency with the solid's constitutive data is a
separate matter, settled by choosing the equilibrium stress $S_e$ of the
reference as the stress found here).

\subsection{Referential form}

To parameterise the shape, pull the state back to a fixed reference body
through an equilibrium mapping $\varphi_e$, with $F_e = D\varphi_e$,
$J_e = \det F_e$. The referential density is $\rho = J_e\,\varrho\circ\varphi_e$,
the first Piola--Kirchhoff stress $P = J_e (T\circ\varphi_e)F_e^{-T}$ and
the referential potential $\zeta = \Phi\circ\varphi_e$. The Piola
identity $J_e(\div \mathbf{a})\circ\varphi_e = \Div\bigl(J_e F_e^{-1}
\mathbf{a}\circ\varphi_e\bigr)$ applied row-wise to $T$ gives
\begin{equation}
  \Div P = J_e (\Div T)\circ\varphi_e
         = \rho\,F_e^{-T}\nabla\zeta ,
  \qquad P\bN = 0 \ \text{on } \partial B,
  \qquad \jump{P\bN} = 0 ,
  \label{eq:equilibrium-ref}
\end{equation}
the traction per referential area $P\bN$ pairing with the Nanson
normal $\bnu = \cof(F_e)\bN$ as in \texttt{doc/slip\_interface.tex}.
The Poisson equation pulls back to the weak form with the tensor
$a_e = J_e F_e^{-1}F_e^{-T}$ (Maitra \& Al-Attar 2024, eq.~21;
\texttt{doc/gravitating\_elasticity.md} \S3), and its source is
$\int_B \rho\,\chi\dx$, \emph{independent of the map} because $\rho$ is
mass per referential volume. These are the referential equilibrium
equations of Al-Attar \& Woodhouse \cite{AW10} and Maitra \& Al-Attar
\cite{MA24}, and nothing more: the fluid condition is
$\dev\bigl(J_e^{-1} P F_e^{T}\bigr) = 0$ in $\Bf$.

\subsection{The mapping may be taken continuous}
\label{sec:continuity}

\begin{proposition}\label{prop:continuous}
Let a static state be described through a labelling that is smooth on
each of $\Bs$, $\Bf$ but slips tangentially across $\Sigma$:
$\varphi_{\mathrm{f}}|_{\Sigma} = \varphi_{\mathrm{s}}|_{\Sigma}\circ\sigma$
with $\sigma$ a diffeomorphism of $\Sigma$ isotopic to the identity.
Then the same physical state is described by data on the same reference
body with a continuous mapping, obtained by relabelling the fluid side
only.
\end{proposition}
\begin{proof}
Physical positions are continuous across a fluid--solid boundary (no
cavitation, no overlap), so $\varphi_{\mathrm{f}}(\Sigma) =
\varphi_{\mathrm{s}}(\Sigma)$ and $\sigma$ is well defined. Because
$\sigma$ is isotopic to the identity, $\sigma^{-1}$ extends (through a
collar of $\Sigma$ in $\Bf$, tapering the isotopy to the identity) to a
diffeomorphism $\xi$ of $\Bf$ onto itself with $\xi|_\Sigma =
\sigma^{-1}$ and $\xi = \mathrm{id}$ away from the collar. Then
$(\varphi_{\mathrm{f}}\circ\xi)|_\Sigma = \varphi_{\mathrm{s}}|_\Sigma$:
the relabelled fluid map matches the solid map on $\Sigma$. Transform the
fluid data by the relabelling laws (Al-Attar \& Crawford \cite{AC16};
\texttt{doc/gravitating\_elasticity.md} \S1), $\tilde\rho = J_\xi\,\rho\circ\xi$
and correspondingly for the constitutive data; the physical fields
$\varrho$, $T$, $\Phi$ are unchanged, hence so is the state.
\end{proof}

The proposition is a statement about \emph{description}, and it holds
because the constitutive data are \emph{defined} at the natural
reference and transported with whatever labelling is chosen. If instead
the data were pinned to a fixed reference while the map varied, the map
would carry physics, and tangential slip would not be removable. That is
precisely the situation of the \emph{perturbation} layer: there the data
are fixed on the reference, a linearised tangential slip is in general
not a relabelling (\texttt{doc/gravitating\_elasticity.md} \S5), and the
slipping-interface formulation is required. For describing and
parameterising an equilibrium, fluid-slip terms are not needed, and the
rest of this note works with continuous $\varphi_e$.

% ======================================================================
\section{Feasibility functionals}
\label{sec:functional}

Fix $(\rho, \varphi_e)$ and the self-consistent $\Phi$. Until
\S\ref{sec:stage2} take $\varphi_e = \mathrm{id}$ (so $\rho = \varrho$);
the mapped forms change nothing structural and are recorded where
needed. The idea is to find the equilibrium stress field that comes
closest to having no deviatoric part in the fluid, and to measure the
remainder. How ``closest'' is measured matters, and three choices are
compared.

\subsection{Admissible stress fields}

For a domain $\Omega \in \{B, \Bf\}$ and a closed subspace
$V \subseteq H^{1}(\Omega)^{3}$ of test fields, define the
\emph{equilibrium-compatible} stress fields
\begin{equation}
  \mathcal{S}(\Omega, V) = \Bigl\{ T \in L^{2}(\Omega; \mathrm{Sym}) :
  \int_{\Omega} T : \eps(\bv)\dV + \int_{\Omega} \bff\cdot\bv\dV = 0
  \ \ \forall\, \bv \in V \Bigr\},
  \label{eq:admissible}
\end{equation}
with $\eps(\bv) = \sym\nabla\bv$. Integrating by parts region by region,
$T \in \mathcal{S}(\Omega, V)$ means $\Div T = \bff$ in each region of
$\Omega$, together with whatever traction conditions the test space
encodes: on parts of $\partial\Omega$ where $V$ is unconstrained,
$T\bn = 0$; across internal interfaces where $V$ is continuous,
$\jump{T\bn} = 0$; on parts of $\partial\Omega$ where $V$ vanishes, no
condition at all. Two instances:
\begin{itemize}
\item \emph{Whole body}: $\Omega = B$, $V = H^{1}(B)^{3}$. This is
  \eqref{eq:equilibrium} without the fluid condition: traction-free
  surface, continuous tractions. The constraint is consistent iff $\bff$
  exerts no net force or torque, $\int_B \bff\cdot\br = 0$ for every
  rigid motion $\br$; for self-gravitation this holds identically,
  because $\int_B \varrho\nabla\Phi\dV$ and $\int_B
  \bx\times\varrho\nabla\Phi\dV$ are double integrals of kernels
  antisymmetric under exchange of the two points.
\item \emph{Fluid alone}: $\Omega = \Bf$, $V = V_{\mathrm{F}} :=
  \{\bv \in H^{1}(\Bf)^{3} : \bv = 0 \text{ on } \Sigma\}$. This is the
  balance $\Div T = \bff$ in the fluid with $T\bn = 0$ on any free fluid
  surface and \emph{no condition on the fluid--solid interface traction}:
  the solid is to absorb whatever traction results.
\end{itemize}

\subsection{The naive functional, and why it fails}
\label{sec:naive}

The first candidate takes the whole-body minimum-deviatoric stress field
of Al-Attar \& Woodhouse \cite[\S3.4]{AW10}, which the library already
generates (\texttt{MinimumDeviatoricEquilibriumStress}), and reads off
its fluid part:
\begin{equation}
  \Ja = \tfrac12\int_{\Bf} |\dev T_m|^{2}\dV , \qquad
  T_m = \operatorname*{arg\,min}_{T \in \mathcal{S}(B, H^1)}
        \tfrac12\int_{B} |\dev T|^{2}\dV .
  \label{eq:Jnaive}
\end{equation}
This is \emph{not} a feasibility certificate: $\Ja$ is in general
positive at models that possess a valid static state.

\begin{proposition}[Leakage]\label{prop:leakage}
Suppose a valid static state exists, and let $T'$ be the valid field
(\,$\dev T' = 0$ in $\Bf$\,) of least solid deviatoric norm
$\int_{\Bs}|\dev T|^2$. Then $\Ja = 0$ if and only if $T'$ is also the
whole-body minimiser $T_m$; this holds if and only if there is a
divergence-free $\bu \in H^1(B)^3$ with $\dev T' = \eps(\bu)$ in $B$ and
$\bu$ rigid on each connected fluid region. Otherwise $\Ja > 0$.
\end{proposition}
\begin{proof}
The objective $Q(T) = \frac12\int_B|\dev T|^2$ is strictly convex in
$\dev T$, so $\dev T_m$ is unique; likewise $\dev T'$ is unique within
the closed affine subset $\mathcal{S}' = \{T \in \mathcal{S}(B,H^1):
\dev T = 0 \text{ on } \Bf\}$. If $\dev T_m = 0$ on $\Bf$ then
$T_m \in \mathcal{S}'$, so $Q(T_m) \ge Q(T')$, and minimality of $T_m$
gives $\dev T_m = \dev T'$. Conversely if $\dev T' = \dev T_m$ then
$\Ja = 0$. Hence $\Ja = 0$ iff $T'$ is a global minimiser. By the
optimality conditions of the whole-body problem (\S\ref{sec:dual} with
uniform weight), $T$ is the global minimiser iff $\dev T = \eps(\bu)$
for some divergence-free $\bu \in H^1(B)^3$; for $T = T'$ this forces
$\eps(\bu) = 0$ on $\Bf$, i.e.\ $\bu$ rigid on each fluid region.
\end{proof}

The mechanism is first-order versus second-order. Perturb $T'$ by a
self-equilibrated field $S$ (\,$\Div S = 0$, continuous tractions,
traction-free surface\,):
\begin{equation}
  Q(T' + tS) = Q(T')
  + t\int_{\Bs}\dev T' : \dev S \dV
  + \tfrac{t^{2}}{2}\int_{B}|\dev S|^{2}\dV .
\end{equation}
The fluid contributes nothing at first order, because $\dev T' = 0$
there; the solid term decreases at first order along any $S$ with
$\int_{\Bs}\dev T':\dev S < 0$. Unless every such $S$ is fluid-shear-free
(the non-generic condition of the proposition), the whole-body minimiser
lowers the solid's deviatoric stress by exporting some into the fluid,
paying only quadratically for it. An aspherical solid enclosing a
barotropic fluid in general must carry deviatoric stress, and the
naive minimiser leaks part of it into the fluid. In a spherically
symmetric hydrostatic model $\dev T' = 0$ everywhere and there is
nothing to leak, which is why spherical tests cannot detect the defect.
The same mechanism is the reason the body's stress field is not
recovered by whole-body minimisation, weighted or not
(\S\ref{sec:recovery}): the minimiser places part of the solid's
unavoidable deviatoric stress in the fluid, where a valid state has
none.

The whole-body value $\Jh := Q(T_m)$ is meaningful in its own right:
$\Jh = 0$ iff the \emph{whole} body admits a hydrostatic state, which is
the defining property of a \emph{hydrostatic figure}
(\S\ref{sec:zeroset}). It is the functional for that problem, not for
equilibrium states with load-bearing solids.

\subsection{The exact fluid-only functional}
\label{sec:JF}

Minimise the fluid's deviatoric stress over fluid stress fields alone,
leaving the interface traction free:
\begin{equation}
  \JF = \inf_{T \in \mathcal{S}(\Bf, V_{\mathrm{F}})}
        \int_{\Bf}\frac{1}{4\mu_{\mathrm{f}}}\,|\dev T|^{2}\dV ,
  \label{eq:JF}
\end{equation}
with a constant reference weight $\mu_{\mathrm{f}} > 0$ (the choice
$\mu_{\mathrm{f}} = \tfrac12$ gives $\tfrac12\int|\dev T|^2$ and matches
the default weight of the library's generator). The weight is written
as a viscosity because that is what it becomes in \S\ref{sec:dual}.

\begin{proposition}[Completion]\label{prop:completion}
Let $\Bs$ be connected. If $T_{\mathrm{f}} \in \mathcal{S}(\Bf,
V_{\mathrm{F}})$, then there is a solid stress $T_{\mathrm{s}}$ with
$\Div T_{\mathrm{s}} = \bff$ in $\Bs$, $T_{\mathrm{s}}\bn =
T_{\mathrm{f}}\bn$ on $\Sigma$ and $T_{\mathrm{s}}\bn = 0$ on
$\partial B \cap \partial\Bs$; the pair is in $\mathcal{S}(B, H^1)$.
\end{proposition}
\begin{proof}
The traction problem for the solid is solvable iff the loads exert no
net force or torque on it (the standard compatibility condition; it is
also \cite[eq.~58]{AW10}, and a solution is supplied e.g.\ by the
minimum-norm generator). Let $\bn_{\mathrm{f}}$ be the outward normal of
$\Bf$. On $\Sigma$ the solid's outward normal is $-\bn_{\mathrm{f}}$, so
the prescribed solid traction there is $-T_{\mathrm{f}}\bn_{\mathrm{f}}$,
and $\Div T_{\mathrm{s}} = \bff$ requires $\int_{\partial\Bs}
T_{\mathrm{s}}\bn = \int_{\Bs}\bff$. The left side is
\begin{equation}
  -\int_{\Sigma} T_{\mathrm{f}}\bn_{\mathrm{f}}\dS
  = -\int_{\partial\Bf} T_{\mathrm{f}}\bn_{\mathrm{f}}\dS
  = -\int_{\Bf}\bff\dV ,
\end{equation}
using that $T_{\mathrm{f}}\bn_{\mathrm{f}} = 0$ on any free fluid surface
and the divergence theorem in $\Bf$. The requirement is therefore
$\int_B\bff\dV = 0$, the vanishing of the total self-gravitational
force. The torque identity is the same computation
with $\bx\times$ inserted, closed by the vanishing of the total
self-gravitational torque (and, for the divergence theorem step, by the
symmetry of $T_{\mathrm{f}}$).
\end{proof}

So the fluid field's forced resultant is exactly what the solid's
equilibrium requires, and nothing is lost by optimising over the fluid
alone. With \emph{several} solid components --- an inner core as well
as a mantle --- the computation above only balances their sum: each
additional solid component imposes its own six force and torque
conditions, which are not automatic. The exact certificate then
enlarges $V_{\mathrm{F}}$ by the traces of the rigid motions of each
additional component (\S\ref{sec:dual}, Remark~\ref{rem:innercore}).
Below, $\Bs$ is connected unless stated.

\begin{proposition}\label{prop:JFexact}
With $\Bs$ connected, $\JF = 0$ if and only if the model possesses a
valid static state \eqref{eq:equilibrium}.
\end{proposition}
\begin{proof}
If a valid state exists, its fluid part is in $\mathcal{S}(\Bf,
V_{\mathrm{F}})$ with zero objective. Conversely, the infimum in
\eqref{eq:JF} is attained (\S\ref{sec:dual}); if it is zero the
minimiser has $\dev T = 0$ in $\Bf$, and Proposition~\ref{prop:completion}
completes it to a valid state.
\end{proof}

\subsection{The weighted whole-body functional}
\label{sec:Jmu}

The third candidate repairs the naive functional by weighting: with a
region-wise positive field $\mu$, equal to $\mu_{\mathrm{f}}$ in the fluid
and $\mu_{\mathrm{s}} \gg \mu_{\mathrm{f}}$ in the solid,
\begin{equation}
  \Jmu = \min_{T \in \mathcal{S}(B, H^1)}
         \int_{B}\frac{1}{4\mu}\,|\dev T|^{2}\dV .
  \label{eq:Jmu}
\end{equation}
This is the library's whole-body generator with a two-valued weight
(its \texttt{mu} coefficient).

\begin{proposition}\label{prop:Jmu}
(i) $\Jmu$ is non-increasing in $\mu_{\mathrm{s}}$, and $\Jmu \ge J_\infty$
where $J_\infty := \inf_{T\in\mathcal{S}(B,H^1)}
\int_{\Bf}|\dev T|^2/(4\mu_{\mathrm{f}})$; with $\Bs$ connected,
$J_\infty = \JF$.
(ii) $\Jmu \le J_\infty + C/\mu_{\mathrm{s}}$, with
$C = \tfrac14\int_{\Bs}|\dev T^{\star}|^{2}$ for any minimiser $T^\star$
of $J_\infty$; in particular at a feasible model
$\Jmu = O(1/\mu_{\mathrm{s}})$.
(iii) $\Jmu \to J_\infty$ as $\mu_{\mathrm{s}} \to \infty$, and the
minimisers converge to the minimiser of $J_\infty$ whose solid
deviatoric norm is least; in particular the solid stress
$\dev T = 2\mu_{\mathrm{s}}\eps(\bu)$ (\S\ref{sec:dual}) stays bounded
and converges, while $\eps(\bu) = O(1/\mu_{\mathrm{s}})$ in the solid.
\end{proposition}
\begin{proof}
(i) The integrand is pointwise non-increasing in $\mu_{\mathrm{s}}$ and
bounded below by its fluid part. $J_\infty = \JF$ for connected $\Bs$:
restriction maps $\mathcal{S}(B,H^1)$ into $\mathcal{S}(\Bf,V_{\mathrm{F}})$,
and Proposition~\ref{prop:completion} shows it is onto.
(ii) Evaluate the objective at $T^\star$.
(iii) From (i)--(ii). For the minimisers: write $\Jmu = \min_T
[\mathcal{F}(T) + \mu_{\mathrm{s}}^{-1}\mathcal{G}(T)]$ with
$\mathcal{F}$ the fluid and $\mathcal{G}$ the (scaled) solid term, both
convex quadratics in $\dev T$. The family is bounded in $L^2$ by (ii); a
weak limit point $\bar T$ minimises $\mathcal{F}$, and comparing with any
minimiser $T^\star$ of $\mathcal{F}$,
$\mathcal{F}(T_\mu) + \mu_{\mathrm{s}}^{-1}\mathcal{G}(T_\mu) \le
\mathcal{F}(T^\star) + \mu_{\mathrm{s}}^{-1}\mathcal{G}(T^\star)$ with
$\mathcal{F}(T_\mu) \ge \mathcal{F}(T^\star)$ gives
$\mathcal{G}(T_\mu) \le \mathcal{G}(T^\star)$, so $\bar T$ minimises
$\mathcal{G}$ among minimisers of $\mathcal{F}$ (the standard
hierarchical limit of a vanishing penalty); strict convexity in the
deviatoric part makes the limit unique, and the convergence strong.
\end{proof}

So $\Jmu$ is a certificate with a floor: at feasible models it is
$O(\norm{\dev T_{\mathrm{s}}}^{2}/\mu_{\mathrm{s}})$, and the floor can
be removed by Richardson extrapolation in $1/\mu_{\mathrm{s}}$ from two or
three contrasts, $\Jmu \approx J_\infty + c_1/\mu_{\mathrm{s}}$. Its
advantages are that it uses the existing generator unchanged, that it
remains exact with any number of solid components (its limit $J_\infty$
contains the force and torque balance of every component, see
Remark~\ref{rem:innercore}), and that its dual (\S\ref{sec:dual}) is the
literal creeping flow of a body with that viscosity structure. Its cost
is conditioning: the Stokes operator with viscosity contrast
$\mu_{\mathrm{s}}/\mu_{\mathrm{f}}$ degrades the preconditioner (a
block-diagonal preconditioner with a viscosity-weighted pressure mass
matrix is the minimum requirement), so the contrast cannot be taken
arbitrarily large.

$\Jmu$ serves as a certificate and as an upper-bound family for $\JF$
(the cross-check of \S\ref{sec:verification}); its minimiser is not a
recovery of the body's stress field. At any finite contrast the fluid
part of the minimiser still carries the residual leak bounded in (ii),
and removing it requires exactly the contrast that conditions the
saddle. The limit field of (iii) is obtained exactly, with no contrast
in any operator, by the construction of \S\ref{sec:recovery}.

\subsection{The zero sets}
\label{sec:zeroset}

\begin{proposition}\label{prop:barotropic}
$\JF = 0$ iff there is $p \in H^{1}(\Bf)$ with
\begin{equation}
  \nabla p = -\rho\nabla\Phi \ \text{in } \Bf,
  \qquad p = 0 \ \text{on free fluid surfaces}.
  \label{eq:hydrostatic}
\end{equation}
At such a state, on each connected component of each level set of
$\Phi$ in $\Bf$ on which $\nabla\Phi \ne 0$, $p$ is constant, and so
is $\rho$ wherever $\rho$ is differentiable; locally $\rho = R(\Phi)$
and $p = P(\Phi)$ with $P' = -R$ (the barotropic structure). A free
fluid surface adjacent to fluid of positive density is a level surface
of $\Phi$. $\Jh = 0$ iff \eqref{eq:hydrostatic} holds on all of $B$ with
$p \in H^1(B)$; then every interface across which $\rho$ jumps is a
level surface of $\Phi$.
\end{proposition}
\begin{proof}
$\dev T = 0$ means $T = -p\,1$; membership of $\mathcal{S}(\Bf,
V_{\mathrm{F}})$ is then the weak form of \eqref{eq:hydrostatic}, with
$p\in H^1$ since $\nabla p = -\rho\nabla\Phi \in L^2$. The tangential
derivative of $p$ along a level set vanishes, so $p$ is constant on each
connected component; taking the curl, $\nabla\rho\times\nabla\Phi = 0$
wherever $\rho$ is differentiable, so $\rho$ is constant there too; by
the implicit function theorem $\rho$ and $p$ are locally functions of
$\Phi$, and \eqref{eq:hydrostatic} gives $P' = -R$. On a free surface
$p = 0$ is constant while $\nabla p = -\rho\nabla\Phi \ne 0$ points
along the normal, so the surface is a level set of $p$, hence locally of
$\Phi$. For $\Jh$, $p \in H^1(B)$ is continuous across interfaces, so
its tangential derivative is too; with
$\nabla_\Sigma p = -\rho\nabla_\Sigma\Phi$ on both sides and $\Phi$
continuously differentiable, $(\rho_+ - \rho_-)\nabla_\Sigma\Phi = 0$:
where the density jumps, the interface is an equipotential.
\end{proof}

The proposition fixes what is and is not an obstruction, and corrects a
tempting over-statement. \emph{Distinct components of one level set are
not forced to carry equal density}: each carries its own branch
$P_k(\Phi)$, and the branches are tied only by single-valuedness of $p$
where components merge and by boundary conditions (a U-tube with two
columns of different density but free surfaces at different heights is
in equilibrium; capped by a solid, the heights need not even differ).
For a fluid region enclosed by a connected solid there is no
obstruction at all: a uniform density $\rho = \text{const}$ satisfies
\eqref{eq:hydrostatic} with $p = -\rho\Phi + \text{const}$, whatever the
geometry, so $\JF = 0$ is always attainable by density control and the
shape of the fluid--solid interface is \emph{not} constrained --- the
solid supports any interface topography. Genuine infeasibility of the
equilibrium-state problem comes from: (a) boundaries that the fluid must
support itself (free fluid surfaces, prescribed density jumps inside the
fluid), which must be equipotentials; (b) the force and torque balance of
each additional solid component; (c) restrictions on the admissible
density (bounds, a fixed parametric form, or the fixed distribution
function of the rearrangement route of \S\ref{sec:energy}). For the
hydrostatic-figure functional $\Jh$, by contrast, every density
interface --- including the core--mantle boundary --- must be an
equipotential, and a fixed aspherical geometry that violates this has
$\Jh > 0$ at the density minimiser, quantifying the infeasibility.

At $\JF = 0$ the minimiser over densities is not unique: every
barotropic density compatible with the constraints is a zero. A prior
selects within the zero set (\S\ref{sec:sobolev}); the rearrangement
route selects differently (\S\ref{sec:energy}).

\begin{remark}[The solid's field]\label{rem:shell}
$\Phi$ is the total potential, $\Phi = \Phi[\rho|_{\Bf}] +
\Phi[\rho|_{\Bs}]$ by linearity, and with the solid fixed (stage~1) the
second part is static data; only the first responds to density updates.
If the solid is a spherically symmetric shell enclosing the fluid, the
shell theorem gives $\nabla\Phi[\rho|_{\Bs}] = 0$ in the cavity, so the
fluid's balance involves the fluid's own self-gravity alone. An
aspherical solid does act on the fluid --- at fixed aspherical geometry
already in stage~1, and through its shape in stage~2.
\end{remark}

% ======================================================================
\section{Duality: the Stokes structure}
\label{sec:dual}

Both functionals are instances of one template. Let $\Omega$, $V$ be as
in \eqref{eq:admissible}, $\mu > 0$ a viscosity field on $\Omega$, and
\begin{equation}
  J = \inf_{T \in \mathcal{S}(\Omega, V)} \int_\Omega \frac{1}{4\mu}|\dev T|^2\dV .
  \label{eq:template}
\end{equation}
$\JF$ is $(\Omega, V, \mu) = (\Bf, V_{\mathrm{F}}, \mu_{\mathrm{f}})$;
$\Jmu$ is $(B, H^1, \mu)$; $\Jh$ is $(B, H^1, \tfrac12)$.

\subsection{Lagrangian and dual}

Introduce a multiplier $\bu \in V$ for the constraint and the Lagrangian
\begin{equation}
  L(T, \bu) = \int_\Omega \frac{1}{4\mu}|\dev T|^2
  - \int_\Omega T : \eps(\bu) - \int_\Omega \bff\cdot\bu .
\end{equation}
$\sup_{\bu} L(T,\bu)$ is the objective if $T \in \mathcal{S}(\Omega,V)$
and $+\infty$ otherwise, so $J = \inf_T\sup_{\bu}L$. For fixed $\bu$,
minimise over $T$ pointwise. Split $T = \dev T + \tfrac13(\tr T)1$:
the trace part enters only through $-\tfrac13\tr T\,\div\bu$, so the
infimum is $-\infty$ unless $\div\bu = 0$; the deviatoric part is
minimised at
\begin{equation}
  \dev T = 2\mu\,\eps(\bu)
  \qquad(\text{using } \dev\eps(\bu) = \eps(\bu) \text{ when } \div\bu = 0),
  \label{eq:devT}
\end{equation}
with value $\int\mu|\eps(\bu)|^2 - 2\int\mu|\eps(\bu)|^2 =
-\int\mu|\eps(\bu)|^2$. Hence the dual problem is
\begin{equation}
  \sup_{\bu \in U} g(\bu), \qquad
  g(\bu) = -\tfrac12 a(\bu,\bu) - \ip{\bff}{\bu}, \qquad
  a(\bu,\bv) = \int_\Omega 2\mu\,\eps(\bu):\eps(\bv)\dV,
  \label{eq:dual}
\end{equation}
over the divergence-free test fields $U = \{\bu \in V : \div\bu = 0\}$,
with $\ip{\bff}{\bu} = \int_\Omega\bff\cdot\bu$. Its maximiser satisfies
\begin{equation}
  a(\bu,\bv) = -\ip{\bff}{\bv} \qquad \forall\,\bv \in U ,
  \label{eq:stokes-weak}
\end{equation}
and at the maximiser, taking $\bv = \bu$, $g(\bu) = \tfrac12 a(\bu,\bu)$.

\begin{proposition}[Strong duality]\label{prop:duality}
Assume $\ip{\bff}{\br} = 0$ for every $\br \in U$ with $a(\br,\br) = 0$.
Then \eqref{eq:stokes-weak} has a solution $\bu$, unique up to that
kernel, there is a pressure $p \in L^2(\Omega)$ with
\begin{equation}
  a(\bu,\bv) - \int_\Omega p\,\div\bv = -\ip{\bff}{\bv}
  \quad\forall\,\bv\in V, \qquad
  \int_\Omega q\,\div\bu = 0 \quad\forall\, q,
  \label{eq:stokes-mixed}
\end{equation}
the field $T_{\min} = -p\,1 + 2\mu\,\eps(\bu)$ attains the infimum
\eqref{eq:template}, and
\begin{equation}
  J = \tfrac12 a(\bu,\bu) = \int_\Omega \mu\,|\eps(\bu)|^2\dV .
  \label{eq:value}
\end{equation}
\end{proposition}
\begin{proof}
On $U$ modulo the kernel, $a$ is coercive by Korn's inequality, and the
compatibility assumption makes the right-hand side a functional on the
quotient; Lax--Milgram gives $\bu$. The pressure exists by the inf--sup
condition for the divergence on $V$ (Girault \& Raviart \cite{GR86}),
which is where divergence-free extensions of boundary data are needed.
By \eqref{eq:stokes-mixed}, $\int T_{\min}:\eps(\bv) = a(\bu,\bv) -
\int p\div\bv = -\ip{\bff}{\bv}$ for all $\bv\in V$, so
$T_{\min}\in\mathcal{S}(\Omega,V)$, and its objective is
$\int\mu|\eps(\bu)|^2 = \tfrac12a(\bu,\bu) = g(\bu)$. Weak duality
($\inf\sup \ge \sup\inf$) then makes both optimal.
\end{proof}

\subsection{Reading the dual}

Equation \eqref{eq:stokes-mixed} is the steady incompressible Stokes
problem
\begin{equation}
  -\nabla p + \Div\bigl(2\mu\,\eps(\bu)\bigr) = \bff = \rho\nabla\Phi ,
  \qquad \div\bu = 0 ,
  \label{eq:stokes-strong}
\end{equation}
i.e.\ $\bu$ is the creeping flow, with viscosity $\mu$, driven by the
gravitational body force $-\rho\nabla\Phi$ against which no hydrostatic
pressure can balance --- the unbalanced buoyancy. The stress
$T_{\min} = -p1 + 2\mu\eps(\bu)$ is its viscous stress. The value
\eqref{eq:value} is $\int\mu|\eps|^2$, which is Rayleigh's dissipation
function of that flow, i.e.\ \emph{half} its rate of viscous
dissipation $\int 2\mu|\eps(\bu)|^2$. (This factor reappears in
\S\ref{sec:energy}.) The boundary conditions are those encoded by $V$:
\begin{itemize}
\item $\JF$: $\bu = 0$ on $\Sigma$ (no slip against the solid) and
  traction-free on free fluid surfaces. The boundary condition is forced
  by the free interface traction: varying $T$ with $\delta T\bn$
  arbitrary on $\Sigma$ in $\int_{\Bf}\delta T:\eps(\bu) = -\int_{\Bf}
  \Div\delta T\cdot\bu + \int_{\partial\Bf}\delta T\bn\cdot\bu$ leaves the
  boundary term, whose vanishing for all $\delta T\bn$ is $\bu|_\Sigma = 0$.
  For a fluid region enclosed by solid, Korn's first inequality makes $a$
  coercive on $U$: \emph{there is no rigid kernel}. The pressure is
  determined only up to an additive constant in each enclosed fluid
  region, since $\int\div\bv = \int_{\partial}\bv\cdot\bn = 0$ for every
  $\bv\in V_{\mathrm{F}}$ there; physically the absolute pressure in a
  sealed container is not fixed by equilibrium. $J$ does not depend on it,
  and it is fixed by a mean-zero constraint per enclosed region. (The
  body's stress field does depend on it; its selection is part of
  \S\ref{sec:recovery}.)
\item $\Jmu$ and $\Jh$: traction-free on $\partial B$, $\bu$ continuous
  across interfaces, $\div\bu = 0$ in the solid as well. The kernel of
  $a$ on $U$ is the global rigid motions (\,$\eps(\bu)=0$ on each region
  plus continuity\,), compatible by the vanishing total force and torque;
  it is projected exactly as in the existing generator. No pressure
  constant is free (traction boundary).
\end{itemize}
In the rigid limit $\mu_{\mathrm{s}}\to\infty$ of $\Jmu$ the solid
velocity becomes rigid on each solid component: with a connected solid,
the limit flow (in the gauge where the solid is at rest) is the no-slip
flow of $\JF$, consistent with Proposition~\ref{prop:Jmu}.

\begin{remark}[Further solid components]\label{rem:innercore}
For an inner core $\Bs^{(2)}$ inside a fluid shell inside a mantle
$\Bs^{(1)}$, the exact fluid-only certificate uses
$V_{\mathrm{F}}^{+} = \{\bv\in H^1(\Bf)^3 : \bv = 0 \text{ on }
\partial\Bs^{(1)},\ \bv = \br \text{ on } \partial\Bs^{(2)} \text{ for some
rigid } \br\}$: six extra unknowns, whose equations in
\eqref{eq:stokes-mixed} are the force and torque balance on the inner
core (inner-core loads $\int_{\Bs^{(2)}}\bff\cdot\br$ included). A rigid
motion has zero flux through the closed inner-core boundary, so the
pressure constant of the fluid shell remains free. This
is the $\mu_{\mathrm{s}}\to\infty$ limit of $\Jmu$ with the mantle at
rest.
\end{remark}

\subsection{Discretisation and inexact solves}

The saddle \eqref{eq:stokes-mixed} is discretised with Taylor--Hood
pairs \cite{TaylorHood}, the pressure one order below the velocity, as
in \texttt{MinimumDeviatoricEquilibriumStress}; $\Jmu$ is exactly that
class with a two-valued \texttt{mu}, and $\JF$ is the same saddle on the
fluid SubMesh with essential (no-slip) conditions on $\Sigma$ and the
per-region pressure constant removed. All statements of this section hold
verbatim for the discrete problem with $U_h$ the discretely
divergence-free subspace, which is what makes the derivatives of
\S\ref{sec:derivative} exact for the discrete functional.

When $\bu$ is only an approximation $\tilde\bu \in U$ with error $\bm{e}$,
the two expressions of the value behave differently:
$g(\tilde\bu) = J - \tfrac12a(\bm{e},\bm{e})$ (a lower bound, error
quadratic in $\bm{e}$, by expanding the quadratic $g$ about its
maximiser), while $\tfrac12a(\tilde\bu,\tilde\bu) = J + a(\bu,\bm{e}) +
O(\bm{e}^2)$ has a first-order error. The value is therefore evaluated as
$g(\tilde\bu)$. The derivative of \S\ref{sec:derivative} is linear in
$\bu$ and carries a first-order error; the saddle tolerance must be set
accordingly.

\subsection{Recovering the body's stress field}
\label{sec:recovery}

The certificate measures infeasibility; the stress field of the whole
body is a separate output. It is recovered exactly by two kinds of solve:
one for the fluid, which is the certificate itself, and one per
connected solid component --- two solves for a fluid shell inside a
mantle, three with an inner core.

\paragraph{Why not a whole-body minimiser.} At uniform weight the
whole-body minimiser leaks the solid's unavoidable deviatoric stress into
the fluid (Proposition~\ref{prop:leakage}). The weighted minimiser of
\eqref{eq:Jmu} reduces the leak to $O(1/\mu_{\mathrm{s}})$
(Proposition~\ref{prop:Jmu}(ii)) but never removes it, and the contrast
that would remove it conditions the saddle (\S\ref{sec:Jmu}). Neither is
used. The construction below delivers the $\mu_{\mathrm{s}}\to\infty$
limit of Proposition~\ref{prop:Jmu}(iii) directly: the fluid part is the
certificate's minimiser, and the solid part minimises the solid's
deviatoric norm given that fluid --- the hierarchical selection of the
limit.

\paragraph{The fluid.} The fluid's stress is the certificate's own
minimiser,
\begin{equation}
  T = -p\,1 + 2\mu_{\mathrm{f}}\,\eps(\bu) \quad\text{in } \Bf ,
\end{equation}
from the fluid-only generator (Proposition~\ref{prop:duality}), with the
rigid border of Remark~\ref{rem:innercore} when there is an inner core.
At feasibility $\dev T\to0$ there and $p$ is the hydrostatic pressure of
\eqref{eq:hydrostatic}, determined up to one constant per enclosed fluid
component; the solver projects the constant mode, so $p$ is returned in
the mean-zero gauge.

\paragraph{Each solid component.} Let $\Omega_{\mathrm{s}}$ be a
connected component of $\Bs$, $\bn$ its outward normal,
$\Sigma_{\mathrm{s}} = \partial\Omega_{\mathrm{s}}\cap\Sigma$ its fluid
interfaces and $\Gamma_{\mathrm{s}} = \partial\Omega_{\mathrm{s}}\cap
\partial B$ its share of the free surface. Its stress is the
minimum-deviatoric field of the component alone, loaded by the fluid's
pressure:
\begin{equation}
  \min\ \tfrac12\int_{\Omega_{\mathrm{s}}}|\dev T|^{2}\dV
  \quad\text{subject to}\quad
  \Div T = \rho\nabla\Phi \ \text{in } \Omega_{\mathrm{s}}, \quad
  T\bn = -p\,\bn \ \text{on } \Sigma_{\mathrm{s}}, \quad
  T\bn = 0 \ \text{on } \Gamma_{\mathrm{s}} .
  \label{eq:solid-recovery}
\end{equation}
This is \texttt{MinimumDeviatoricEquilibriumStress}
(\texttt{background.hpp}) on the component's SubMesh, with its optional
interface-pressure load, which is applied against the outward normal of
the component regardless of the orientation the boundary elements
inherit from the parent mesh. The load is the pressure part of the fluid
traction; the viscous part $2\mu_{\mathrm{f}}\eps(\bu)\bn$ is
$O(J^{1/2})$ and vanishes at feasibility. Solving per component gives
each component its own rigid-mode projection, which is required: with an
inner core the solid is disconnected, and each component's loads must
balance on their own. At a restored state they do. On the mantle this is
the computation of Proposition~\ref{prop:completion}; on the inner core
it is exactly the content of the rigid-border certificate
(Remark~\ref{rem:innercore}), whose six extra equations are the core's
force and torque balance. Away from feasibility the rigid-mode
projection removes the residual imbalance, which is of the order of the
fluid's residual.

\paragraph{The pressure gauge.} The constant $c$ of an enclosed fluid's
pressure is a genuine parameter of the equilibrium family: $p\mapsto
p + c$ leaves the fluid's balance and its deviatoric stress unchanged and
adds the load $-c\,\bn$ on $\Sigma_{\mathrm{s}}$. Its effect on the solid
depends on how much of $\partial\Omega_{\mathrm{s}}$ the fluid covers.
\begin{itemize}
\item \emph{Entirely fluid-loaded boundary} (the inner core,
  $\partial\Omega_{\mathrm{s}} = \Sigma_{\mathrm{s}}$). A uniform pressure
  on a closed surface exerts no net force or torque, and $T - c\,1$ solves
  \eqref{eq:solid-recovery} with $p + c$: the gauge shifts the stress
  isotropically and $\dev T$ is invariant.
\item \emph{Partly free boundary} (the mantle: fluid interface plus free
  surface). A constant interface pressure against a traction-free outer
  surface is not carried isotropically; its response is the thick-shell
  Lam\'e solution, which has genuine deviatoric stress. The gauge
  matters.
\end{itemize}
The generator's own principle resolves it. Let $T_0$ solve
\eqref{eq:solid-recovery} with the gauged pressure, and $L$ with zero body
force and unit interface pressure ($L\bn = -\bn$ on $\Sigma_{\mathrm{s}}$,
$L\bn = 0$ on $\Gamma_{\mathrm{s}}$). The generator is linear in its data
and its deviatoric output is unique, so $T(c) = T_0 + c\,L$ is the
equilibrium family in $c$. Minimising $\norm{\dev T(c)}^{2}$ over $c$,
\begin{equation}
  c^{\star} = -\frac{(\dev T_0, \dev L)}{\norm{\dev L}^{2}} ,
  \label{eq:gauge-recovery}
\end{equation}
with $(\cdot,\cdot)$ the $L^2$ deviatoric pairing over
$\Omega_{\mathrm{s}}$: one extra solve. Minimising over $T$ at fixed $c$
and then over $c$ is the joint minimum, so $T(c^\star)$ is the
component's minimum-deviatoric field with the gauge free. The recovered
$c^\star$ is the physical pressure datum the free surface imposes: in a
hydrostatic configuration it is the constant that makes the interface
pressure continuous with the solid's lithostatic pressure, which
vanishes at the free surface. With several enclosed fluid components, or
several partly free components sharing one fluid, the constants form a
vector and \eqref{eq:gauge-recovery} becomes a small least-squares
problem; in the inner core, fluid shell and mantle configuration only
the mantle sees the constant, and the scalar formula is the whole of it.

\paragraph{A verification fact.} On a spherically symmetric configuration
with a barotropic fluid the valid state has $\dev T = 0$ everywhere, and
the recovery finds it: every component's deviatoric norm sits at the
discretisation floor (any $c \ne c^\star$ would leave the mantle with the
Lam\'e shear $|c - c^\star|\,\norm{\dev L}$). The fluid's norm is exactly
$(2J)^{1/2}$ at $\mu_{\mathrm{f}} = \tfrac12$, by the value-function
identity \eqref{eq:value}: $J = \int_{\Bf}|\dev T|^2/(4\mu_{\mathrm{f}})
= \tfrac12\norm{\dev T}^2_{L^2(\Bf)}$. The whole-body map of
\texttt{examples/equilibrium\_density.cpp} is assembled from these
pieces, on the common scale of \S\ref{sec:stopping}.

% ======================================================================
\section{The derivative with respect to density}
\label{sec:derivative}

Stage~1: the shape is fixed and the control is the fluid density. All
three functionals depend on $\rho$ only through the load $\bff =
\rho\nabla\Phi[\rho]$; the spaces $V$, $U$ and the form $a$ are
independent of $\rho$.

\subsection{The Poisson problem and its closure}
\label{sec:poisson}

Let $D \supset B$ be the ball bounded by the DtN sphere, of radius $b$;
the region $D\setminus B$ is the buffer. Outside $D$ the potential is
harmonic and decaying, $\Phi = \sum_{lm}\Phi_{lm}(b/r)^{l+1}Y_{lm}$, so on
$\partial D$, $\partial_r\Phi = -\Lambda\Phi$ with $\Lambda$ the
Dirichlet-to-Neumann operator, acting on degree $l$ by $(l+1)/b$:
symmetric and positive semi-definite. Multiplying
\eqref{eq:poisson-strong} by $\chi\in H^1(D)$ and integrating by parts
over $D$,
\begin{equation}
  b(\Phi,\chi) := \frac{1}{4\pi G}\Bigl[\int_D \nabla\Phi\cdot\nabla\chi\dV
  + \int_{\partial D}\chi\,\Lambda\Phi\dS\Bigr]
  = -\int_B \rho\,\chi\dV
  \qquad\forall\,\chi\in H^1(D) .
  \label{eq:poisson-weak}
\end{equation}
The form $b$ is symmetric and, in three dimensions, coercive on $H^1(D)$.
This is the library's Poisson machinery (Laplacian plus DtN, with the DtN
truncated at some degree; the truncated form is still symmetric, which
is all that is used below). Write $\Phi[\rho]$ for the solution operator;
it is linear.

\begin{remark}[Two dimensions]\label{rem:2d}
The 2-D DtN form has no degree-zero term (the exterior harmonic of
degree zero is the logarithm, which does not decay), so the
Laplace--DtN form $b$ is singular in the constant: the potential is
determined up to a constant, and $b(\cdot,\chi) = \ell(\chi)$ is solvable
only if $\ell(1) = 0$. The forward load violates this whenever the body
has net mass $m = \int_B\rho\dV$: the exterior monopole is logarithmic,
its flux through $\partial D$ is not zero, and a closure without a
degree-zero term cannot carry it. The load is made compatible by
subtracting the mass as a uniform flux through the outer boundary,
\begin{equation}
  b(\Phi,\chi) = -\int_B\rho\,\chi\dV
  + \frac{m}{|\partial D|}\int_{\partial D}\chi\dS ,
\end{equation}
which is exactly the monopole's flux (uniform on the circle), so the
interior gradient is unchanged and only the constant is left open.
Solutions --- the potential, its perturbations and the adjoint --- are
taken in the zero-boundary-mean gauge $\int_{\partial D}\Phi\dS = 0$.
The adjoint load $\int\rho\bu\cdot\nabla\chi$ vanishes at $\chi = 1$ and
needs no correction. In this gauge the envelope formulas below hold
discretely without correction. In any other gauge they do not: the
perturbed forward problem carries the flux term of $\delta m =
\int_B\delta\rho$, and the pairing of the adjoint against a direction of
mass $\delta m$ acquires the extra term
\begin{equation}
  \frac{1}{|\partial D|}\Bigl(\int_{\partial D} w\dS\Bigr)\,\delta m .
\end{equation}
It vanishes on mass-preserving directions, but not in the dual vector,
which the loops pair with arbitrary directions before the mass
projection and in Hessian products. $J$ itself reads only $\nabla\Phi$
and is gauge-free. (The Poisson solver of \texttt{equilibrium\_figures.hpp}
projects the constant and returns the zero-boundary-mean gauge.)
\end{remark}

\subsection{Differentiating the value function}

\begin{proposition}\label{prop:envelope}
Let $\ell_\rho(\bv) := -\ip{\rho\nabla\Phi[\rho]}{\bv}$ and $\bu$ the
solution of \eqref{eq:stokes-weak}. Then
\begin{equation}
  \delta J = \delta\ell_\rho(\bu)
  = -\int_\Omega \bigl(\delta\rho\,\nabla\Phi + \rho\,\nabla\delta\Phi\bigr)\cdot\bu\dV ,
  \qquad \delta\Phi = \Phi[\delta\rho] .
  \label{eq:dJ1}
\end{equation}
\end{proposition}
\begin{proof}
Let $\mathsf{S}: U' \to U$ be the solution operator of
\eqref{eq:stokes-weak} on the quotient by the kernel: $a(\mathsf{S}\ell,\bv)
= \ell(\bv)$. It is self-adjoint, $\ell_1(\mathsf{S}\ell_2) =
a(\mathsf{S}\ell_1,\mathsf{S}\ell_2)$. Then $J = \tfrac12 a(\mathsf{S}\ell_\rho,
\mathsf{S}\ell_\rho) = \tfrac12\ell_\rho(\mathsf{S}\ell_\rho)$, and
$\delta J = \tfrac12\delta\ell_\rho(\mathsf{S}\ell_\rho) +
\tfrac12\ell_\rho(\mathsf{S}\delta\ell_\rho) = \delta\ell_\rho(\bu)$
by self-adjointness. The kernel is harmless: $\ell_\rho(\br) = 0$ for
every rigid $\br$ and every $\rho$ (vanishing total force and torque),
so also $\delta\ell_\rho(\br) = 0$.
\end{proof}

This is the envelope theorem (Danskin \cite{Danskin}) for the value
function $J(\rho) = \max_{U} g(\cdot\,;\rho)$: the derivative is the
partial derivative of the Lagrangian at the saddle point, and no
transposed saddle solve is needed --- \emph{the primal multiplier $\bu$
is the adjoint variable}, because the saddle is self-adjoint and the
objective is its own value. The proof shows this directly for stage~1;
the Lagrangian form of the envelope theorem is what carries over to
stage~2, where $U$ itself depends on the control. For $\JF$ the pairing
in \eqref{eq:dJ1} is over $\Bf$ only; for $\Jmu$ it is over $B$, the
solid's slow flow included.

\subsection{The Poisson adjoint}

The first term of \eqref{eq:dJ1} is already local in $\delta\rho$. The
second, $T_2 := \int_\Omega\rho\nabla\delta\Phi\cdot\bu$, is not. Define
the adjoint potential $w\in H^1(D)$ by
\begin{equation}
  b(\chi, w) = \int_\Omega \rho\,\bu\cdot\nabla\chi\dV
  \qquad\forall\,\chi\in H^1(D),
  \label{eq:adjoint}
\end{equation}
the same symmetric Poisson operator with a different load. Then, taking
$\chi = \delta\Phi$ in \eqref{eq:adjoint} and $\chi = w$ in the
perturbed \eqref{eq:poisson-weak} (\,$b(\delta\Phi,\chi) =
-\int\delta\rho\,\chi$\,),
\begin{equation}
  T_2 = b(\delta\Phi, w) = -\int_B \delta\rho\, w\dV .
\end{equation}
With \eqref{eq:dJ1}, for perturbations supported in $\Bf$,
\begin{equation}
  \boxed{\;\delta J = \int_{\Bf}\delta\rho\,\bigl(w - \bu\cdot\nabla\Phi\bigr)\dV,
  \qquad J'(\rho) = w - \bu\cdot\nabla\Phi \ \ \text{in } \Bf\; }
  \label{eq:gradient}
\end{equation}
as an $L^2$ density. The formula holds unchanged for perturbations of
the solid density (with $\bu$ the solid velocity, zero for $\JF$), which
is used in the verification of \S\ref{sec:verification}.

\paragraph{Strong form and interpretation.} Integrating
\eqref{eq:adjoint} by parts, $w$ solves $\nabla^2 w = 4\pi G\,\div(\rho\bu)$
in the distributional sense (with $\rho\bu$ extended by zero, so the
jumps of $\rho\bu\cdot\bn$ act as surface sources) and the same DtN
closure. Let $\phi^{1}[\bu]$ be the Eulerian potential perturbation
produced by transporting the density with displacement $\bu$: the
Eulerian density change is $-\div(\rho\bu)$, so $\nabla^2\phi^1 =
-4\pi G\div(\rho\bu)$. Hence $w = -\phi^1[\bu]$ and
\begin{equation}
  J'(\rho) = -\bigl(\phi^{1}[\bu] + \bu\cdot\nabla\Phi\bigr)
  = -\,\Delta_{\bu}\Phi ,
  \label{eq:lagrangian}
\end{equation}
minus the \emph{Lagrangian} perturbation of the potential along $\bu$ ---
in the language of \texttt{doc/gravitating\_elasticity.md} \S3, minus the
referential potential perturbation $\zeta^1 = \phi^1 + \bu\cdot\nabla\Phi$
generated by $\bu$. An independent check of the sign by Green's
functions: with $\Phi[\sigma](\bx) = -G\int\sigma(\mathbf{y})|\bx-\mathbf{y}|^{-1}$,
$T_2 = -\int\div(\rho\bu)\,\delta\Phi = \int\delta\rho\,\phi^1[\bu]$ by
the symmetry of the kernel, and $-T_2 = \int\delta\rho\,w$ agrees.

\paragraph{Mapped form.} With a fixed non-identity $\varphi_e$ (stage~1
on an aspherical reference) the same steps give, with $\zeta =
\Phi\circ\varphi_e$, the pulled-back velocity $\bU = F_e^{-1}\bu$ and the
mapped Poisson form $b_e$ (Laplacian with $a_e = J_eF_e^{-1}F_e^{-T}$, DtN
unchanged because $\varphi_e = \mathrm{id}$ near $\partial D$):
$b_e(\chi,w) = \int\rho\,\bU\cdot\nabla\chi$ and $J' = w - \bU\cdot\nabla\zeta$.

\subsection{Discrete exactness}

For the discrete functional $J_h(\rho_h)$, defined by the discrete
Poisson form $b_h$, the discrete load $\ell_h(\bv_h) =
-\int\rho_h\nabla\Phi_h\cdot\bv_h$ and the Taylor--Hood saddle, every step
above holds in finite dimensions. The derivative vector
$j_i = \partial J_h/\partial\rho_i = \int_{\Bf}\psi_i\,(w_h -
\bu_h\cdot\nabla\Phi_h)$, with $w_h$ from the discrete adjoint
$b_h(\chi_h,w_h) = \int\rho_h\bu_h\cdot\nabla\chi_h$, is then the
\emph{exact} gradient of $J_h$ (up to solver tolerances), provided the
same quadrature is used for the load in the forward and adjoint
assemblies. Finite-difference checks therefore converge to round-off,
not to discretisation level.

\subsection{Homogeneity and an identity}

Without rotation, $\ell_\rho$ is a quadratic form in $\rho$ (the
product of $\rho$ with the linear $\nabla\Phi[\rho]$), $\bu$ is linear in
$\ell_\rho$, and $J$ is quadratic in $\bu$: $J$ is a homogeneous quartic
in the total density $\rho$ on $B$, $J(t\rho) = t^4J(\rho)$. Euler's
identity gives
\begin{equation}
  \int_B \rho\,J'(\rho)\dV = 4J ,
  \label{eq:euler}
\end{equation}
a check that exercises the whole chain. Directly:
$\int\rho\,\phi^1[\bu] = \int(-\div\rho\bu)\Phi = \int\rho\bu\cdot\nabla\Phi$
by the mutual-energy symmetry, so $\int\rho J' = -2\ip{\bff}{\bu} =
2a(\bu,\bu) = 4J$. (The integral is over the whole body: $J'$ on the
solid is $w - \bu\cdot\nabla\Phi$ as well.)

Second derivatives --- Hessian actions with respect to density and
shape, their Gauss--Newton part, and Newton--Krylov methods --- are the
subject of \S\ref{sec:second}.

% ======================================================================
\section{Energy interpretation and the advection route}
\label{sec:energy}

\subsection{The functional as a squared gradient}

Let $E = \tfrac12\int_B\rho\,\Phi[\rho] = -\frac{1}{8\pi G}\int_{\R^3}
|\nabla\Phi|^2 \le 0$ be the gravitational energy (the identity by
\eqref{eq:poisson-strong} and integration by parts). Move the material
with a velocity $\bv$ that is a \emph{rearrangement}: an element of $U$
(divergence-free in the fluid, and for $\JF$ zero on $\Sigma$; for
$\Jmu$ divergence-free throughout). The density is transported,
$\partial_t\rho = -\div(\rho\bv)$, and since $E$ is a symmetric
quadratic form,
\begin{equation}
  \frac{\mathrm{d}E}{\mathrm{d}t} = \int\partial_t\rho\,\Phi
  = -\int\div(\rho\bv)\,\Phi = \int\rho\,\bv\cdot\nabla\Phi
  = \ip{\bff}{\bv} ,
  \label{eq:dE}
\end{equation}
(boundary terms vanish because $\rho\bv$ is extended by zero and $\Phi$
is continuous). So $\bv\mapsto\ip{\bff}{\bv}$ is the derivative of $E$
along rearrangements, and by \eqref{eq:stokes-weak}
\begin{equation}
  \bu = -\operatorname{grad}_a E ,
  \qquad J = \tfrac12\norm{\operatorname{grad}_a E}_a^{2} ,
\end{equation}
where $\operatorname{grad}_a$ is the Riesz representative in the
\emph{dissipation metric} $a$ on $U$. Minimising $J$ is thus a
least-squares formulation of the first-order optimality conditions of
$E$ on the space of rearrangements; the stationarity conditions of $E$
under rearrangement are exactly the equilibrium conditions, with the
pressure as the multiplier of incompressibility (this is the classical
variational characterisation of equilibria of incompressible
self-gravitating fluids).

Consequently $J = 0$ at \emph{every} critical point of $E$, stable or
not. The second variation along the flow $X_t$ of a fixed $\bv\in U$,
$\rho_t = \rho\circ X_t^{-1}$, is
\begin{equation}
  \frac{\mathrm{d}^2E}{\mathrm{d}t^2}\Big|_{0}
  = \int\ddot\rho\,\Phi + \int\dot\rho\,\Phi[\dot\rho],
  \qquad \dot\rho = -\bv\cdot\nabla\rho, \quad
  \ddot\rho = \bv\cdot\nabla(\bv\cdot\nabla\rho).
\end{equation}
Integrating the first term by parts ($\div\bv = 0$, $\bv\cdot\bn = 0$ on
the fluid boundary) and using $\rho = R(\Phi)$ at a barotropic state,
\begin{equation}
  \frac{\mathrm{d}^2E}{\mathrm{d}t^2}\Big|_{0}
  = -\int_{\Bf} R'(\Phi)\,(\bv\cdot\nabla\Phi)^2\dV
    \;-\; \frac{1}{4\pi G}\int_{\R^3}\bigl|\nabla\Phi[\bv\cdot\nabla\rho]\bigr|^2\dV .
\end{equation}
The second (self-gravity) term is non-positive but smoothing: for $\bv$
localised at scale $\ell$ it is smaller than the first by $O(\ell^2)$.
Hence $R' \le 0$ (density non-increasing with $\Phi$, i.e.\ outwards) is
necessary for stability, and an unstably stratified barotropic state
($R' > 0$ somewhere) is a saddle of $E$ with $J = 0$. This is the
stability criterion of an incompressibly rearranged fluid: parcels
carry their density unchanged; a compressible fluid would be judged
against the adiabatic gradient instead (the Brunt--V\"ais\"al\"a
criterion), which the rearrangement model does not represent.

\subsection{Gradient flow: the advection route}

The alternative outer loop transports the density along $\bu$,
\begin{equation}
  \partial_t\rho + \bu\cdot\nabla\rho = 0 \quad\text{in } \Bf ,
  \label{eq:advection}
\end{equation}
recomputing $\Phi$ and $\bu$ as $\rho$ changes. By \eqref{eq:dE} and
\eqref{eq:stokes-weak},
\begin{equation}
  \frac{\mathrm{d}E}{\mathrm{d}t} = \ip{\bff}{\bu} = -a(\bu,\bu) = -2J :
\end{equation}
this is the gradient flow of $E$ in the dissipation metric --- the
quasi-static viscous relaxation of an incompressible fluid of viscosity
$\mu$ --- and the energy decays at twice the functional's value (the
full dissipation rate). For $\Jmu$ it is literally the relaxation of the
whole body with the two-valued viscosity, the solid creeping at
$O(1/\mu_{\mathrm{s}})$ (a stage-1 implementation freezes the geometry,
consistently to that order); for $\JF$ the solid is held at rest.

Properties, each immediate from \eqref{eq:advection} with
$\div\bu = 0$ and $\bu\cdot\bn = 0$ on the fluid boundary:
\begin{itemize}
\item The flow map is volume-preserving, so mass is conserved and the
  \emph{distribution function} $|\{\rho > c\}\cap\Omega_k|$ of every
  connected fluid region $\Omega_k$ is invariant. Advection answers
  ``rearrange \emph{this} material into equilibrium'' --- for instance,
  rearrange the material of a reference Earth model's core into a state
  consistent with an aspherical mantle --- whereas the optimisation route
  answers ``find \emph{a} nearby equilibrium density''.
\item Mass cannot move between disconnected fluid regions, and the
  per-region distribution is fixed: this is the advection route's form of
  the density restrictions of \S\ref{sec:zeroset}, item (c).
\item $E$ decreases strictly until $J = 0$; generically the flow leaves
  saddles (unstable stratifications) and selects stable states, whereas
  minimising $J$ converges to whatever zero is nearest, stable or not,
  and needs a prior for selection.
\item For a fluid region enclosed by a fixed solid, $-E$ restricted to
  the region is (up to a linear term from the solid's field) the strictly
  convex, weakly sequentially continuous functional
  $\rho\mapsto\frac{1}{8\pi G}\norm{\nabla\Phi[\rho]}^2$; by Burton's
  theory of convex maximisation over rearrangement classes
  \cite{Burton1987} its maximum over the weak closure of the
  rearrangement class of the initial density is attained \emph{within}
  the class. A minimiser of $E$ among rearrangements of the given
  material therefore exists without mixing. In the spherically symmetric
  case it is explicit: by the shell theorem (Remark~\ref{rem:shell}) the
  solid's term is constant on the class, and by the Riesz rearrangement
  inequality (Lieb \& Loss \cite[Thm.~3.7]{LiebLoss}) the self-energy is
  minimised by the symmetric decreasing rearrangement $\rho^{*}$ of the
  initial density about the centre.
\end{itemize}

\subsection{Cost}

Per step, both loops need one Poisson solve for $\Phi$ and one saddle
solve for $\bu$; the optimisation route needs one more Poisson solve (the
adjoint $w$) plus the Riesz map of \S\ref{sec:sobolev}, the advection
route one transport step. Transport of $\rho$ needs care with the
discrete divergence: Taylor--Hood velocities are only discretely
divergence-free against the pressure space, so the conservative form
$\partial_t\rho + \div(\rho\bu_h) = 0$ conserves mass exactly but not the
distribution function, and the advective form the reverse; an
exactly divergence-free reconstruction of $\bu_h$ (an $H(\div)$
post-processing) restores both.

\subsection{The discrete advection step: a caution}
\label{sec:advection-caution}

The advection route looks like the safest loop: it is the physical
relaxation, it needs no line search on $J$, and its energy decreases at
the known rate $2J$. A direct Eulerian discretisation nevertheless fails,
in a way worth understanding, because the failure is structural rather
than a matter of tuning.

\paragraph{Acceptance on $J$ rejects the true flow.} $J$ need not
decrease along the gradient flow of $E$. One dimension shows why: for
$f(x) = \cos x$ started near the maximum $x = 0$, the gradient flow
descends $f$ while $\tfrac12 f'(x)^2 = \tfrac12\sin^2x$ \emph{grows} until
the inflection point. In general $J = \tfrac12\norm{\operatorname{grad}_aE}_a^2$
changes along the flow at the rate $-\ip{\operatorname{grad}_aE}
{\operatorname{Hess}_aE\,\operatorname{grad}_aE}_a$, up to terms
proportional to $\operatorname{grad}_aE$ itself. So $J$ grows wherever the
flow runs along negative curvature of $E$ --- leaving an unstable
arrangement, or overturning. A step test on $J$ therefore rejects
legitimate steps of the flow. The correct test is on $E$, Armijo-style,
using the exact rate $\mathrm{d}E/\mathrm{d}t = -2J$.

\paragraph{Acceptance on $E$ admits collapse.} The implemented Eulerian
step is the weak form of $\partial_t\rho = -\div(\rho\bu)$,
\begin{equation}
  (\delta\rho_h, v) = \Delta t\int_{\Bf}\rho_h\,\bu_h\cdot\nabla v\dV
  \qquad\forall\,v ,
  \label{eq:weak-advection}
\end{equation}
which moves the derivative onto the test function (the control is never
differentiated) and conserves mass exactly ($v = 1$; $\bu_h = 0$ on
$\Sigma$). It is not a rearrangement, for two reasons.
\begin{itemize}
\item $\bu_h$ is divergence-free only weakly, against the pressure space.
  Integrating \eqref{eq:weak-advection} by parts exposes the source
  $-\rho\,\div\bu_h$: a compression, not a transport.
\item Even for an exactly divergence-free $\bu$, the explicit step is not
  range-preserving. With $\delta = -\Delta t\,\Pi(\bu\cdot\nabla\rho)$,
  $\Pi$ the $L^2$ projection, $(\delta,\rho) = \Delta t\int\bu\cdot\nabla
  (\rho^2/2) = -\Delta t\int\tfrac12\rho^2\div\bu = 0$. Hence
  $\norm{\rho + \delta}^2 = \norm{\rho}^2 + \norm{\delta}^2$: each step
  raises the amplitude at $O(\Delta t^2)$. Mass is conserved; the
  distribution function is not.
\end{itemize}
Off the rearrangement manifold $E$ is unbounded below at fixed mass:
concentrating mass toward the centre sends $\norm{\nabla\Phi}^2$ to
infinity. Along any ray $E$ is a \emph{concave} quadratic,
\begin{equation}
  E(\rho + t\delta) = E(\rho) + t\int\delta\,\Phi
  - \frac{t^2}{8\pi G}\norm{\nabla\Phi[\delta]}^2 ,
\end{equation}
so an energy-decrease test with step growth rewards long steps, and long
explicit steps are exactly where the off-manifold amplitude lives. The
step control thus selects the amplitude channel, and the loop runs into
\emph{discrete gravitational collapse}: $E$ falls, $J$ rises, density
piles up. Capping the step by a CFL condition slows the drift but cannot
stop it. CFL governs the stability of the transport discretisation, and
this is a consistency defect --- a leak out of the constraint manifold
--- which the objective rewards.

\paragraph{As implemented, and the structural fix.} The implemented loop
accepts on $E$ (Armijo with $\mathrm{d}E/\mathrm{d}t = -2J$), caps the
step by CFL, and enforces an absolute ceiling on $J$ (twice its starting
value). With the ceiling it stalls rather than collapses, which is the
honest outcome. The structural fix is a range-preserving step: a
semi-Lagrangian composition $\rho^{n+1} = \rho^{n}\circ X_n^{-1}$ along
the discrete flow map, or better $\rho^n = \rho^0\circ(X_n\circ\cdots\circ
X_1)^{-1}$ with the composed map carried as the unknown. Values are then
resampled, never combined, so bounds and (up to the map's discrete volume
error) the distribution function are preserved. The composed map is a
relabelling-type map, which connects the advection route to the mapping
layer.

% ======================================================================
\section{Derivatives are not gradients: Sobolev metrics}
\label{sec:sobolev}

\subsection{Derivative, metric, gradient}

$J'(\rho)$ is a dual element: discretely, the vector $j$ with
$j_i = \partial J_h/\partial\rho_i$. A gradient needs an inner product,
represented by a symmetric positive-definite Gram matrix $G$: $g =
G^{-1}j$, the unique vector with $g^{T}G\,\delta = j^{T}\delta$ for all
$\delta$. The choice of $G$ is the implicit preconditioner of every
first-order method. The $L^2$ choice $G = M$ (the mass matrix) gives
$g = M^{-1}j$, the $L^2$ projection of the density \eqref{eq:gradient}.
It is a reasonable start here because $J'$ is built from smooth fields
($\bu\in H^1$, $\phi^1 \in H^1$, $\nabla\Phi$ bounded), with two caveats:
discretely, $\nabla\Phi_h$ jumps across element faces, so $M^{-1}j$
carries mesh-scale oscillations that are not damped and accumulate over
iterations; and the Stokes flow has singular gradients at re-entrant
edges and where free fluid surfaces meet solid, imprinting local
structure on $J'$.

\subsection{The $H^1$ metric}

On a domain $X$ (\S\ref{sec:extension}) take
\begin{equation}
  \ip{\rho_1}{\rho_2}_{1} = \int_X\bigl(\alpha\,\rho_1\rho_2
  + \beta\,\nabla\rho_1\cdot\nabla\rho_2\bigr)\dV ,
  \qquad G_1 = \tilde A := \alpha M + \beta K ,
\end{equation}
with $K$ the stiffness matrix. The Riesz solve is the Helmholtz-type
problem $(\alpha - \beta\nabla^2)g = J'$; its Green's function in three
dimensions is $e^{-r/L}/(4\pi\beta r)$ with smoothing length
$L = \sqrt{\beta/\alpha}$: the gradient is the derivative averaged over
balls of radius $\sim L$, and components of $J'$ at wavelengths below $L$
are suppressed by $(1 + \beta k^2/\alpha)^{-1}$.

\subsection{Regularity and the $H^2$ default}

In three dimensions $H^1 \hookrightarrow L^6$ only; $H^s
\hookrightarrow C^0$ requires $s > 3/2$. A Riesz representative in $H^s$
of an arbitrary dual element of $H^s$ is in $H^s$, so a metric of order
$s > 3/2$ guarantees continuous updates whatever the regularity of the
derivative (discretisation noise, or future parameter fields with
rougher derivatives). The working default is $s = 2$, realised by
iterating the $H^1$ Gram. With $\mathcal{A}_h = M^{-1}\tilde A$ the
discrete operator $\alpha - \beta\nabla^2_h$ (self-adjoint in the
$M$-inner product), the order-$2$ inner product is
$\ip{\mathcal{A}_h\rho_1}{\mathcal{A}_h\rho_2}_{L^2}$, i.e.
\begin{equation}
  G_2 = \tilde A M^{-1}\tilde A , \qquad
  g = \tilde A^{-1} M \tilde A^{-1} j :
\end{equation}
two solves with the same operator and one mass product. The caveat: the
composition carries the boundary condition of $\tilde A$ through both
applications, so the norm is $\norm{(\alpha - \beta\nabla^2_X)\rho}_{L^2}$
with the condition built into the operator domain --- spectrally
$(\alpha + \beta(-\nabla^2_X))^{2}$, not the condition-free $H^2(X)$
norm. This is harmless because the metric is posed on an extension
domain, whose boundary is kept several smoothing lengths away from the
control region.

\subsection{Extension domain}
\label{sec:extension}

The parameter is regarded as living on a domain $X \supseteq \Bf$ ---
the whole body, or the ball including the buffer shell --- and the
forward model reads only its restriction to $\Bf$; the derivative $j$ is
supported in $\Bf$, and the Riesz solve spreads the gradient into
$X\setminus\Bf$, where it is inert. This removes any boundary condition
on the parameter at the fluid--solid interface, which is precisely where
the density action concentrates and where an artificial Dirichlet or
Neumann condition would bias every update. The outer condition (Dirichlet
or natural) acts at $\partial X$, at distance $\gg L$ from $\Bf$, and its
influence decays like $e^{-d/L}$. The extended field is continuous: it is
a phantom in the solid, not the solid's density, which is separate data.
A consequence: updates in a metric of order $s \ge 1/2$ cannot create
discontinuities in $\rho|_{\Bf}$ (step functions are in $H^s$ only for
$s < 1/2$); discontinuities present in the initial density are kept,
none are generated.

\subsection{Fractional orders}

For a positive self-adjoint operator $\mathcal{A}$ and $0 < \sigma < 1$,
Balakrishnan's formula \cite{Balakrishnan1960}
\begin{equation}
  \mathcal{A}^{-\sigma} = \frac{\sin\pi\sigma}{\pi}\int_0^\infty
  t^{-\sigma}\,(t + \mathcal{A})^{-1}\,\mathrm{d}t
  \label{eq:balakrishnan}
\end{equation}
follows from the scalar identity $\int_0^\infty t^{-\sigma}(t+\lambda)^{-1}
\mathrm{d}t = \lambda^{-\sigma}\pi/\sin\pi\sigma$ ($\lambda > 0$) and
the spectral theorem. For the metric of order $s = n + \sigma$ the
gradient is $g = \mathcal{A}_h^{-s}M^{-1}j$; using
$(t + M^{-1}\tilde A)^{-1}M^{-1} = (tM + \tilde A)^{-1}$,
\begin{equation}
  \mathcal{A}_h^{-\sigma}M^{-1}j = \frac{\sin\pi\sigma}{\pi}
  \int_0^\infty t^{-\sigma}\,\bigl((\alpha + t)M + \beta K\bigr)^{-1} j\,\mathrm{d}t ,
\end{equation}
each integrand evaluation a Helmholtz solve with a shifted, better
conditioned operator. With $t = e^{y}$ the integrand decays exponentially
in both directions (like $e^{(1-\sigma)y}$ as $y\to-\infty$ and
$e^{-\sigma y}$ as $y\to\infty$), and the trapezoidal (sinc) rule in $y$
converges exponentially in the number of nodes (Bonito \& Pasciak
\cite{BonitoPasciak2015}). The order $s$ becomes a continuous
regularisation parameter, realised with nothing but shifted solves of
existing machinery; the same Riesz and fractional-power infrastructure
serves every parameter field of the library's inverse problems.

\subsection{Constraints and priors in the metric}

\emph{Mass.} The mass functional $m(\rho) = \int_{\Bf}\rho$ has
derivative vector $m_i = \int_{\Bf}\psi_i$, independent of $\rho$; its
Riesz representative $r = G^{-1}m$ is fixed and cached. The
$G$-orthogonal projection of a gradient onto the mass-preserving
directions is
\begin{equation}
  g_\perp = g - \frac{m^{T}g}{m^{T}r}\,r ,
\end{equation}
which satisfies $m^{T}g_\perp = 0$ and $\ip{g_\perp}{r}_G = 0$: one scalar
correction per step. Further linear constraints (centre of mass and
products of inertia in the rotating case, \S\ref{sec:rotation}) add
cached representatives and a small Gram matrix.

\emph{Prior.} A Tikhonov term $\frac{\lambda}{2}(\rho-\rho_{\mathrm{ref}})^T
G(\rho-\rho_{\mathrm{ref}})$ in the \emph{same} metric has derivative
$\lambda G(\rho-\rho_{\mathrm{ref}})$ and hence gradient
$\lambda(\rho-\rho_{\mathrm{ref}})$: plain shrinkage, no extra solve.
This requires $\rho_{\mathrm{ref}}$ on $X$ (a continuous extension of the
fluid reference profile, satisfying the outer condition) and the prior
measured with exactly the Gram used for the gradient.

% ======================================================================
\section{Stage 2: the shape as a control}
\label{sec:stage2}

\subsection{Gauge}

Admit $\varphi_e$ among the controls. The pair $(\rho,\varphi_e)$
carries the relabelling gauge: for a diffeomorphism $\xi$ of $B$
preserving every region and equal to the identity near $\partial D$,
$(J_\xi\,\rho\circ\xi,\ \varphi_e\circ\xi)$ describes the same physical
state. Since $J$ depends only on the physical state, it is
gauge-invariant. Infinitesimally, with $\xi = \mathrm{id} + t\bfeta$
($\bfeta$ tangent to $\partial B$ and to every interface), $J_\xi \approx
1 + t\Div\bfeta$ and $\rho\circ\xi \approx \rho + t\bfeta\cdot\nabla\rho$,
so the gauge directions are
\begin{equation}
  (\delta\rho, \delta\varphi_e) = (\Div(\rho\bfeta),\ F_e\bfeta),
\end{equation}
the referential counterparts of the trivial displacements of
Friedman \& Schutz \cite{FriedmanSchutz1978}. The derivative annihilates
them, so in exact arithmetic the gradient in any metric is orthogonal to
the gauge and descent never drifts along it. Discretely, invariance is
only approximate, and spurious gauge components accumulate over many
iterations --- the optimisation analogue of the semi-convergence of
\texttt{doc/gauge\_penalty\_iteration.tex}. The redundancy is eliminated
rather than projected: $\varphi_e = \varphi[h]$ is parameterised by
interface-shape coefficients $h$ (topography of the fluid--solid
interfaces and of the surface, e.g.\ in spherical harmonics), extended
into the regions by a fixed rule and equal to the identity at the DtN
sphere. If $h\mapsto$ (physical interface shapes) is injective, a
relabelling with $\varphi[h]\circ\xi = \varphi[h']$ forces $h = h'$ and
$\xi = \mathrm{id}$: no gauge remains. This matches the mapping layer and
keeps the computational interfaces spherical.

\subsection{The shape derivative}

The envelope theorem now needs its Lagrangian form, because the
divergence constraint depends on the map. Pulled back, the saddle reads:
$a_\varphi(\bu,\bv) = \int 2\mu\,\eps_\varphi(\bu):\eps_\varphi(\bv)\,J$
with $\eps_\varphi(\bu) = \sym(D\bu\,F^{-1})$; divergence
$\tr(D\bu\,F^{-1})J$; load $\ell(\bv) = -\int\rho\,(F^{-T}\nabla\zeta)
\cdot\bv$; Poisson $b_\varphi$ with $a(F) = JF^{-1}F^{-T}$ and the
map-independent source $-\int\rho\chi$. With
$\mathcal{L}(\bu,p;\rho,\varphi) = -\tfrac12a_\varphi(\bu,\bu) +
\ell(\bu) + \int p\,\tr(D\bu F^{-1})J$ and $(\bu,p)$ the saddle point,
\begin{equation}
  \delta J = -\tfrac12\,\delta a_\varphi(\bu,\bu)
  - \int\rho\,(\delta F^{-T}\nabla\zeta)\cdot\bu
  + \delta b_\varphi(\zeta, w)
  + \int p\,\delta\bigl[\tr(D\bu F^{-1})J\bigr] ,
  \label{eq:shape-derivative}
\end{equation}
with $\delta F = D\delta\varphi$, $\delta(F^{-1}) = -F^{-1}\delta F F^{-1}$,
$\delta J = J\tr(F^{-1}\delta F)$, $w$ the mapped adjoint of
\S\ref{sec:derivative}, and $\delta b_\varphi(\zeta,w) = \frac{1}{4\pi G}
\int\ip{a'(\delta\varphi)\nabla\zeta}{\nabla w}$ with $a'$ as in
\texttt{doc/gravitating\_elasticity.md} \S3.1. The third term follows as
before: differentiating $b_\varphi(\zeta,\chi) = -\int\rho\chi$ gives
$b_\varphi(\delta\zeta,\chi) = -\delta b_\varphi(\zeta,\chi)$, and the
$\nabla\delta\zeta$ part of the load derivative is $-b_\varphi(\delta\zeta,w) =
\delta b_\varphi(\zeta,w)$. All terms are volume integrals: these are the
moving-domain transport terms, the penalised fluid region moving with
$\varphi_e$, written in the distributed (volume) form of the shape
derivative. For $\JF$ the no-slip condition $\bu = 0$ on $\Sigma$ is
map-independent in referential coordinates. The equivalent boundary
(Hadamard) form concentrates the derivative on the interfaces as normal
speed times jumps of energy densities, paired through the Nanson normal
$\bnu = \cof(F_e)\bN$ with the referential measure (the Jacobian
warning of \texttt{doc/slip\_interface.tex}); the volume form is
preferred discretely (Hiptmair, Paganini \& Sargheini
\cite{HiptmairPaganiniSargheini2015}). An inner core's rigid trace
$\bu = W\varphi_e + \mathbf{t}$ on its boundary (Remark~\ref{rem:innercore})
depends on the map and contributes a further interface term
$\int_{\partial\Bs^{(2)}}(P\bN)\cdot W\delta\varphi\dS$ with the fluid
traction per referential area, paired through Nanson; its sign and the
rest of that case are to be fixed with the stage-2 implementation.

The field-level form of \eqref{eq:shape-derivative}, its projection onto
interface parameterisations, the role of the lift and the second-order
(Gauss--Newton) shape block are in \S\ref{sec:shape-second}.

The derivative of the state with respect to interface location is what
the interface-shift perturbation benchmark certifies
(\texttt{benchmarks/perturbation/README.md}); shift-type maps must be
used in their interpolated-$F$ form and with gradient kinks aligned with
mesh attributes, as documented there.

\subsection{What stage 2 is for}

By \S\ref{sec:zeroset}, a fluid region enclosed by a connected solid is
always feasible at fixed shape, so for the equilibrium-state functional
$\JF$ the shape is a genuine control only where fluid supports its own
boundary --- free fluid surfaces (oceans, wholly fluid bodies), fluid
interfaces with prescribed density jumps --- where an inner core must be
balanced, or where the density is restricted (\S\ref{sec:zeroset}(c),
e.g.\ the rearrangement semantics). For the hydrostatic-figure functional
$\Jh$, every density interface must become an equipotential, and the
fixed-aspherical-geometry case is a useful intermediate diagnostic:
$\Jh > 0$ at the stage-1 minimiser quantifies the infeasibility of a
non-equipotential interface, and stage~2 drives it to zero.

\subsection{Rotation}
\label{sec:rotation}

A steady rotation $\bm{\Omega}$ adds the centrifugal potential
$\psi = -\tfrac12|\bm{\Omega}\times\bx|^2$ (house convention: the
centrifugal acceleration is $-\nabla\psi$), $\bff = \rho\nabla(\Phi+\psi)$,
turning equilibrium states into equilibrium figures. $\psi$ is
independent of $\rho$, so \eqref{eq:gradient} acquires only
$-\bu\cdot\nabla\psi$, and $J$ is no longer homogeneous. Compatibility is
no longer automatic: $\int\rho\nabla\psi = -\Omega^2 M\bx_{\perp,\mathrm{cm}}$
vanishes iff the centre of mass lies on the rotation axis, and the
torque $\int\rho\,\bx\times\nabla\psi$ (with $\bm{\Omega}$ along
$\mathbf{e}_3$, components $\Omega^2\int\rho\,x_2x_3$ and
$-\Omega^2\int\rho\,x_1x_3$) vanishes iff the axis is a principal axis of
inertia. Both are linear in $\rho$ and enter as additional projected
constraints (\S\ref{sec:sobolev}). Deferred: Coriolis and the
Tisserand-frame machinery in the perturbation layer, finite-strain
reference changes, coupling to viscoelastic evolution.

% ======================================================================
\section{Second derivatives and second-order adjoints}
\label{sec:second}

\subsection{Why second order is routine here}

In the referential framework every control is a field on the fixed
reference domain: the density $\rho$ on $B$ (or on the extension domain
$X$), and in stage~2 the equilibrium mapping $\varphi_e$ on $D$, or the
interface-shape coefficients $s$ through which it is parameterised.
Every form --- the Stokes saddle, the load, the Poisson operator with
its DtN closure --- is an integral over fixed reference regions, on a
fixed mesh with fixed quadrature points, whose integrand depends on the
controls pointwise through $\rho$ and through $F = D\varphi_e$ (via
$F^{-1}$, $J = \det F$, $\cof F$). There are no boundary perturbations
anywhere. Derivatives of all orders are therefore ordinary Gateaux
derivatives of integrands, and second-order adjoints are obtained by
applying the Lagrangian argument of \S\ref{sec:derivative} once more.

\begin{remark}[Contrast with moving-domain shape calculus]
Classical shape optimisation perturbs the domain itself by vector fields
$V$; the first shape derivative then has the Hadamard structure (it
depends only on $V\cdot\bn$ on the boundary), and the second derivative,
defined by composing two perturbations of the identity, carries
curvature terms and is in general not symmetric: it depends on how the
second perturbation is transported by the first (Delfour \& Zol\'esio
\cite{DelfourZolesio}). Parameterising the shape by a map on a fixed
reference trades all of this for $F$-dependence of the integrands. A
function of the coefficients $s\in\R^{N_s}$ has an ordinary, symmetric
second derivative; normal speeds, curvature and the jump structure on the
interfaces are absorbed into the dependence on $F$, which the mapping
layer keeps explicit and analytic at the quadrature points.
\end{remark}

\paragraph{Discrete setting.} Write the Taylor--Hood saddle of
\eqref{eq:stokes-mixed} as
\begin{equation}
  \mathcal{K}\begin{bmatrix}\bu\\ p\end{bmatrix}
  = \begin{bmatrix}-\mathsf{f}\\ 0\end{bmatrix},
  \qquad
  \mathcal{K} = \begin{bmatrix} \mathsf{A} & -\mathsf{D}^{T}\\
                 -\mathsf{D} & 0\end{bmatrix},
  \label{eq:saddle-discrete}
\end{equation}
with $\mathsf{A}$ the matrix of $a$, $\mathsf{D}$ the divergence pairing
and $\mathsf{f}_i = \int\rho\,\nabla\Phi\cdot\bv_i$ (mapped:
$\int\rho\,F^{-T}\nabla\zeta\cdot\bv_i$) the load vector; and the Poisson
problem as $\mathsf{B}\Phi = -\mathsf{m}(\rho)$ with $\mathsf{B}$ the
Laplacian-plus-DtN matrix (mapped: with $a(F)$) and $\mathsf{m}(\rho)_k =
\int\rho\chi_k$, independent of the map. Then $J = \tfrac12\bu^{T}
\mathsf{A}\bu$ and, by Proposition~\ref{prop:envelope},
\begin{equation}
  \mathrm{d}J[d] = -\bu^{T}\,\partial\mathsf{f}[d] .
  \label{eq:dJ-discrete}
\end{equation}

\subsection{The density Hessian}
\label{sec:hessian}

In stage~1, $\mathcal{K}$ and $\mathsf{B}$ do not depend on $\rho$, and
$\mathsf{f}$ depends on it bilinearly, $\mathsf{f}(\rho) =
\mathsf{f}_2(\rho,\rho)$ with $\mathsf{f}_2(\rho_1,\rho_2)_i =
\int\rho_1\nabla\Phi[\rho_2]\cdot\bv_i$ (a centrifugal term adds a
part linear in $\rho$ and changes nothing below). Hence
\begin{equation}
  \partial\mathsf{f}[d] = \bigl(d\,\nabla\Phi + \rho\,\nabla\delta\Phi_d,\ \bv_i\bigr),
  \qquad
  \partial^2\mathsf{f}[d_1,d_2] = \bigl(d_1\nabla\delta\Phi_{2} +
  d_2\nabla\delta\Phi_{1},\ \bv_i\bigr),
  \qquad \delta\Phi_k = \Phi[d_k],
\end{equation}
each $\delta\Phi$ one Poisson solve. The state sensitivity solves the
\emph{same} saddle:
\begin{equation}
  \mathcal{K}\begin{bmatrix}\delta\bu_k\\ \delta p_k\end{bmatrix}
  = \begin{bmatrix}-\partial\mathsf{f}[d_k]\\ 0\end{bmatrix} .
  \label{eq:sensitivity}
\end{equation}
Differentiating \eqref{eq:dJ-discrete} (with $d = d_1$) in direction
$d_2$,
\begin{equation}
  H[d_1,d_2] = -\delta\bu_2^{T}\,\partial\mathsf{f}[d_1]
               - \bu^{T}\partial^2\mathsf{f}[d_1,d_2] .
\end{equation}
By the first row of \eqref{eq:sensitivity} for $d_1$,
$-\partial\mathsf{f}[d_1] = \mathsf{A}\delta\bu_1 - \mathsf{D}^T\delta p_1$,
and $\mathsf{D}\delta\bu_2 = 0$ by its second row for $d_2$, so the
pressure term drops:
\begin{equation}
  \boxed{\;H[d_1,d_2] = \underbrace{\delta\bu_2^{T}\mathsf{A}\,\delta\bu_1}_{H_{\mathrm{GN}}}
  \;\underbrace{-\;\bu^{T}\partial^2\mathsf{f}[d_1,d_2]}_{H_{\mathrm{R}}}\;}
  \label{eq:hessian-split}
\end{equation}
Both terms are manifestly symmetric in $(d_1,d_2)$.

\paragraph{The Gauss--Newton term.} $J = \tfrac12\norm{r}^2$ with the
residual $r(\rho) = (2\mu)^{1/2}\eps(\bu(\rho)) \in L^2(\Omega;\mathrm{Sym})$,
a least-squares problem, and $H_{\mathrm{GN}}[d_1,d_2] =
\ip{\partial r[d_1]}{\partial r[d_2]}_{L^2} = a(\delta\bu_1,\delta\bu_2)$
is its Gauss--Newton matrix: symmetric positive semi-definite. Its order:
$H_{\mathrm{GN}}[d,d] = \sup_{\bv\in U}\ip{d\nabla\Phi +
\rho\nabla\delta\Phi_d}{\bv}^2/a(\bv,\bv)$ is the squared $U'$-norm of the
force perturbation, an $H^{-1}$-type norm of $d$: $H_{\mathrm{GN}}$ is
smoothing of order $-2$ on densities, hence compact on $L^2$. Its kernel
contains every $d$ whose force perturbation is a gradient --- at $J = 0$,
the tangent directions of the barotropic family.

\paragraph{The residual term.} $H_{\mathrm{R}}$ pairs the residual
velocity $\bu$ with the bilinear part of the load. Since
$\norm{\bu}_a = (2J)^{1/2}$, $|H_{\mathrm{R}}[d,d]| = O(J^{1/2})$: it
vanishes on the zero set, where Gauss--Newton is exact, and away from it
it is indefinite (it is what makes the quartic $J$ non-convex), so the
full Hessian can be indefinite.

\paragraph{Hessian-vector product.} For fixed $d_2 = d$, the dual vector
$h$ with $h^{T}d_1 = H[d_1,d]$ follows by the Poisson-adjoint step of
\S\ref{sec:derivative} applied to each term. The Gauss--Newton term is
\eqref{eq:dJ-discrete} with $\bu$ replaced by $\delta\bu_d$:
$-\int(d_1\nabla\Phi + \rho\nabla\delta\Phi_1)\cdot\delta\bu_d$; the
residual term is $-\int(d_1\nabla\delta\Phi_d + d\,\nabla\delta\Phi_1)
\cdot\bu$. The parts non-local in $d_1$ are $-\int\nabla\delta\Phi_1\cdot
(\rho\,\delta\bu_d + d\,\bu)$, and with one adjoint
\begin{equation}
  b(\chi,\hat w) = \int\bigl(\rho\,\delta\bu_d + d\,\bu\bigr)\cdot\nabla\chi
  \quad\forall\chi
  \qquad\Longrightarrow\qquad
  -\int\nabla\delta\Phi_1\cdot(\rho\,\delta\bu_d + d\,\bu)
  = \int d_1\,\hat w ,
\end{equation}
exactly as for $w$. Hence
\begin{equation}
  H d = \hat w - \delta\bu_d\cdot\nabla\Phi - \bu\cdot\nabla\delta\Phi_d ,
  \label{eq:hessvec}
\end{equation}
where $\hat w = -\delta\phi^1$ is minus the Eulerian potential of
$-\div(\rho\,\delta\bu_d + d\,\bu)$ (in strong form
$\nabla^2\hat w = 4\pi G\div(\rho\,\delta\bu_d + d\,\bu)$, same operator
and closure as $w$): \eqref{eq:hessvec} is the derivative of
\eqref{eq:lagrangian}. The Gauss--Newton product drops $d\,\bu$ from the
adjoint source and the last term.

These are the second-order adjoints of the pair. Because the saddle is
self-adjoint and its value is the objective, the second-order saddle
adjoint coincides with the forward sensitivity $\delta\bu_d$; the
second-order Poisson adjoint is $\hat w$, with the source of $w$
differentiated.

\paragraph{Cost.} One product needs $\delta\Phi_d$ (Poisson),
$\delta\bu_d$ (saddle) and $\hat w$ (Poisson): one saddle and two Poisson
solves, the same solves as one value-and-gradient evaluation ($\Phi$,
$\bu$, $w$), i.e.\ between one and two value evaluations depending on
the cost of a Poisson solve relative to a saddle solve, plus the
assembly of right-hand sides. No operator depends on $\rho$:
$\mathcal{K}$, $\mathsf{B}$ and their preconditioners are built once for
the whole stage-1 run.

\subsection{Relation to the energy}

At a zero of $J$ (a critical point of $E$ on rearrangements,
\S\ref{sec:energy}), take a rearrangement direction $d = -\div(\rho\bv)$,
$\bv \in U$. The force perturbation $\ip{\partial\mathsf{f}[d]}{\bv'}$ is
the $\rho$-derivative of $E'[\bv'] = \ip{\bff}{\bv'}$ along the flow of
$\bv$, i.e.\ the second variation $Q_E(\bv,\bv')$ of $E$ on
rearrangements (at a critical point the term from composing two flows,
$\ip{\bff}{\nabla_{\bv}\bv'}$, vanishes because $\bff$ annihilates $U$).
Hence $\delta\bu_d = -\operatorname{Hess}_aE\,\bv$ and
\begin{equation}
  H[d,d] = H_{\mathrm{GN}}[d,d] = \norm{\operatorname{Hess}_aE\,\bv}_a^{2}
  \qquad\text{at } J = 0 :
\end{equation}
on rearrangement directions the Hessian of $J$ is the \emph{square} of the
second-variation operator of $E$ in the dissipation metric. The second
variation of $E$ therefore lives in the Gauss--Newton term, and squaring
loses its sign: $J$ has a minimum at stable and unstable stratifications
alike, as found in \S\ref{sec:energy}. The sign survives in the
linearisation $-\operatorname{Hess}_aE$ of the advection flow, which is
why that route discriminates stability. The residual term carries no
second-variation information; it couples the current imbalance $\bu =
-\operatorname{grad}_aE$ to the curvature of the load map and is
responsible for the non-convexity of $J$ away from equilibrium.

\subsection{Newton--Krylov for density}
\label{sec:newton-krylov}

Each outer step solves approximately the Gauss--Newton model on the mass
plane,
\begin{equation}
  \min_{\delta:\ \mathsf{m}^{T}\delta = 0}\
  j^{T}\delta + \tfrac12\delta^{T}\bigl(H_{\mathrm{GN}} + \lambda G\bigr)\delta,
\end{equation}
where $\lambda G$ is the exact Hessian of the Tikhonov prior of
\S\ref{sec:sobolev} (its gradient being $\lambda(\rho - \rho_{\mathrm{ref}})$,
added to $j$). The method is truncated conjugate gradients
(Steihaug \cite{Steihaug1983}) with matrix-free products
\eqref{eq:hessvec} and the preconditioner $z = \Pi_G G^{-1}r$, where
$\Pi_G$ is the $G$-orthogonal projection onto the mass plane of
\S\ref{sec:sobolev}. Preconditioning with a projection that maps into the
constraint plane keeps every iterate on it (projected CG; Gould, Hribar
\& Nocedal \cite{GouldHribarNocedal2001}).

\begin{itemize}
\item \emph{Singular directions.} $H_{\mathrm{GN}}$ is exactly singular at
  $J = 0$ along the barotropic family. In stage~1 there is no gauge
  direction (the density alone carries none); the field parameterisation
  $(\rho,\varphi_e)$ of stage~2 would add the gauge directions to the
  kernel at critical points, and the interface-shape parameterisation
  removes them. Rearrangement directions are \emph{not} kernel directions
  in general: by the previous subsection they lie in the kernel only where
  $\operatorname{Hess}_aE$ does (neutral directions such as rearrangements
  within a uniform region, which do not change $\rho$ at all). Truncated
  CG tolerates semi-definiteness; with $\lambda > 0$ the model is
  positive definite; a trust region bounds the step otherwise.
\item \emph{Indefiniteness.} Away from the minimum the full Hessian can
  be indefinite; full Newton then needs negative-curvature handling.
  Gauss--Newton is positive semi-definite by construction and asymptotically
  exact on the zero set.
\item \emph{The preconditioner.} The Riesz map is the natural
  preconditioner in a specific sense: the preconditioned operator
  $G^{-1}(H_{\mathrm{GN}} + \lambda G) = \lambda I + G^{-1}H_{\mathrm{GN}}$
  is the identity plus a compact operator ($H_{\mathrm{GN}}$ has order
  $-2$, $G$ order $2s \ge 0$), so CG converges superlinearly with
  mesh-independent iteration counts. $G^{-1}H_{\mathrm{GN}}$ alone is not
  well conditioned --- it is more compact still --- so without the prior
  the inner CG is an iterative regularisation, the metric deciding that
  smooth components are resolved first and truncation deciding where to
  stop.
\item Bounds are handled as in \S\ref{sec:algorithms} (projection of the
  Newton iterate, with backtracking).
\end{itemize}

\subsection{Shape: the field-level derivative and its projections}
\label{sec:shape-second}

\paragraph{The field dual.} Equation \eqref{eq:shape-derivative} is
linear in $\delta F = D\delta\varphi$. Written as a vector-field dual on
the reference,
\begin{equation}
  \Gamma(\delta\varphi) := \partial_\varphi J[\delta\varphi]
  = \int_{\Omega}\Psi : D\delta\varphi\dx
  + \int_{D}\Psi_{g} : D\delta\varphi\dx ,
  \label{eq:field-dual}
\end{equation}
with ($\eps_\varphi = \sym(D\bu F^{-1})$, $\bm{g} = F^{-T}\nabla\zeta$,
$\bU = F^{-1}\bu$, and $A:B = \tr(A^TB)$)
\begin{align}
  \Psi &= 2\mu J\bigl[F^{-T}D\bu^{T}\eps_\varphi F^{-T}
          - \tfrac12|\eps_\varphi|^{2}F^{-T}\bigr]
          + \rho\,\bm{g}\otimes\bU
          + pJ\bigl[\tr(D\bu F^{-1})F^{-T} - F^{-T}D\bu^{T}F^{-T}\bigr],
  \\
  \Psi_{g} &= \frac{1}{4\pi G}\Bigl[(\nabla w\cdot a\nabla\zeta)F^{-T}
          - F^{-T}\nabla w\otimes a\nabla\zeta
          - F^{-T}\nabla\zeta\otimes a\nabla w\Bigr] ,
\end{align}
the three groups of $\Psi$ coming, in order, from $-\tfrac12\delta
a_\varphi$, from the explicit $F$ in the load, and from the divergence
constraint, and $\Psi_g$ from $\delta b_\varphi$ with $a'$ of
\texttt{doc/gravitating\_elasticity.md} \S3.1 (each obtained from
$\delta F^{-1} = -F^{-1}\delta FF^{-1}$, $\delta J = J\,F^{-T}:\delta F$
and $A:(B\,\delta F\,C) = (B^TAC^T):\delta F$). $\Psi$ and $\Psi_g$ are the
energy--momentum (Eshelby-type) tensors of the Stokes and Poisson
Lagrangians. A centrifugal potential composed with the map, whose load
is $-\int\rho\,(\nabla\psi\circ\varphi_e)\cdot\bu$, adds the term
$-\int\rho\,\bu\cdot(\nabla^2\psi\circ\varphi_e)\,\delta\varphi$, in
$\delta\varphi$ itself rather than its gradient. $\Gamma$ is assembled once per outer
iteration from fields already computed for the density gradient ($\bu$,
$p$, $\zeta$, $w$): \emph{no additional solve}. It is a dual vector on
the finite element space of the map.

\paragraph{Projection onto a parameterisation.} Any shape
parameterisation $\varphi_e = \varphi[s]$ enters only through its
tangent fields $\bm{\theta}_i = \partial\varphi/\partial s_i$: the
interface displacement of coefficient $s_i$ (e.g.\ a spherical harmonic
on one interface) lifted into the volume by a fixed rule --- the mapping
layer's harmonic or tapered extensions --- and vanishing near the DtN
sphere. Then
\begin{equation}
  \frac{\partial J}{\partial s_i} = \Gamma(\bm{\theta}_i) ,
\end{equation}
one dot product per coefficient. One derivation serves every
parameterisation; the interface-harmonic one is a particular pairing.

\paragraph{Lifts and the gauge.} Two lifts $\bm{\theta}, \bm{\theta}'$ of
the same interface motion differ by a field vanishing on every interface,
$\bm{\theta} - \bm{\theta}' = F_e\bfeta$. This is \emph{not} by itself a
gauge direction: the gauge direction is $(\Div(\rho\bfeta), F_e\bfeta)$,
moving the map \emph{and} transporting the referential density.
Gauge invariance of $J$ gives the identity
\begin{equation}
  \Gamma(F_e\bfeta) = -\int_B J'_\rho\,\Div(\rho\bfeta)\dx
  = \int_B \rho\,\nabla J'_\rho\cdot\bfeta\dx
  \label{eq:gauge-identity}
\end{equation}
for every $\bfeta$ tangent to the interfaces and vanishing near
$\partial D$ ($J'_\rho = w - \bU\cdot\nabla\zeta$ on the whole body, as in
\S\ref{sec:derivative}). Hence, at fixed referential density, the
projected derivative depends on the lift by
$\int\rho\nabla J'_\rho\cdot\bfeta$. This is physical, not a defect: a
different lift moves different material when the referential density is
held fixed. The difference vanishes where $\rho = 0$ (lifts differing
only in the buffer give identical derivatives, the extension gauge of
\texttt{doc/gravitating\_elasticity.md} \S3), and where $J'_\rho$ is
constant: in the fluid at a density-stationary point, where $J'_\rho$
equals the mass multiplier. In the solid, whose density is data, the lift
is part of the model's definition of how material moves with the
interfaces. A lift-independent derivative is obtained by pairing each
lift with its gauge-compensating density change, holding the
\emph{Eulerian} density fixed away from the interfaces:
\begin{equation}
  \Gamma_{\mathrm{E}}(\bm{\theta}) := \Gamma(\bm{\theta})
  + \int_B J'_\rho\,\Div_{\mathrm{r}}\bigl(\rho F_e^{-1}\bm{\theta}\bigr)\dx ,
\end{equation}
with $\Div_{\mathrm{r}}$ the region-wise divergence. By
\eqref{eq:gauge-identity}, $\Gamma_{\mathrm{E}}$ depends only on the
interface traces of $\bm{\theta}$. This is the derivative of classical
(Eulerian) shape calculus, Hadamard structure included, recovered from
the referential one. Which pairing is used is a choice of coordinates on
the joint $(\rho, s)$ control space; a joint stationary point is the same
in both. The referential pairing conserves each region's mass
automatically, matching the rearrangement semantics of
\S\ref{sec:energy}.

\paragraph{Discrete gauge diagnostic.} Discretely
\eqref{eq:gauge-identity} holds only approximately. The lift of each
parameter is therefore fixed once and for all, and the defect
\begin{equation}
  \Gamma_h(\bm{\theta}) - \Gamma_h(\bm{\theta}')
  - \int\rho_h\nabla J'_{\rho,h}\cdot\bfeta
  \qquad(\text{equivalently } \Gamma_{\mathrm{E},h}(\bm{\theta}) -
  \Gamma_{\mathrm{E},h}(\bm{\theta}'))
\end{equation}
for two lifts of one interface motion is a built-in measure of the
discrete gauge error. It is the optimisation analogue of the
change-of-variables identity checks of the mapped generators
(\texttt{MappedGeneratorsMatchTransformedMesh}).

\subsection{Shape second order: Gauss--Newton}

With shape controls, $\mathcal{K}$, $\mathsf{f}$ and $\mathsf{B}$ all
depend on $s$ through $F$. Define the residual on the fixed reference,
\begin{equation}
  r(\rho, s) = (2\mu J_\varphi)^{1/2}\,\eps_\varphi(\bu)
  \in L^2(\Omega;\mathrm{Sym}),
  \qquad J = \tfrac12\norm{r}^{2}_{L^2} ,
\end{equation}
(the weight $J_\varphi^{1/2}$ absorbs the volume change, so the space
does not move), and the Gauss--Newton matrix as the Gram matrix of its
first derivative, $H_{\mathrm{GN}}[\cdot,\cdot] =
\ip{\partial r[\cdot]}{\partial r[\cdot]}_{L^2}$. For density directions
this coincides with \eqref{eq:hessian-split}. For a shape direction
$\bm{\theta}$,
\begin{equation}
  \partial r[\bm{\theta}] = (2\mu J_\varphi)^{1/2}\Bigl[
  \tfrac12(F^{-T}:D\bm{\theta})\,\eps_\varphi(\bu)
  - \sym\bigl(D\bu F^{-1}D\bm{\theta}F^{-1}\bigr)
  + \eps_\varphi(\delta\bu_{\bm{\theta}})\Bigr],
\end{equation}
which needs only the \emph{first-order} state sensitivity, obtained by
differentiating \eqref{eq:saddle-discrete}:
\begin{equation}
  \mathcal{K}\begin{bmatrix}\delta\bu_{\bm{\theta}}\\ \delta p_{\bm{\theta}}\end{bmatrix}
  = -\,\partial\mathcal{K}[\bm{\theta}]\begin{bmatrix}\bu\\ p\end{bmatrix}
  - \begin{bmatrix}\partial\mathsf{f}[\bm{\theta}]\\ 0\end{bmatrix},
  \qquad
  \mathsf{B}\,\delta\zeta_{\bm{\theta}} = -\partial\mathsf{B}[\bm{\theta}]\,\zeta ,
\end{equation}
where $\partial\mathsf{f}[\bm{\theta}]$ contains the explicit $F$-dependence
of the load and $\nabla\delta\zeta_{\bm{\theta}}$ (the map-independent
Poisson source contributes nothing). Per shape direction this is one
Poisson and one saddle solve with the \emph{current} operators, whose
preconditioners are reused across all directions within an outer
iteration. The second row is non-zero, $\mathsf{D}\delta\bu_{\bm{\theta}} =
-\partial\mathsf{D}[\bm{\theta}]\bu$: the divergence constraint moves with
the map.

Two points distinguish this from the density case. First, the
Gauss--Newton matrix must be defined through the residual, not by
discarding a term of the value-function expansion. The second-order
envelope formula $\mathrm{d}^2J = \partial^2_{ss}\mathcal{L} -
c^{T}\mathcal{K}_L^{-1}c$, with $\mathcal{K}_L$ the Hessian of the
Lagrangian in $(\bu,p)$ and $c = \partial_{(\bu,p)}\partial_s\mathcal{L}$,
has a sensitivity term that is indefinite as soon as $c$ has a pressure
component, because $\mathcal{K}_L$ is a saddle matrix. For density
directions $c$ has none, and the two definitions agree. Second, full
Newton needs second derivatives of every form with respect to $F$ (and,
for a parameterisation nonlinear in $s$, $\partial^2\varphi/\partial s^2$)
together with second-order adjoint solves. Gauss--Newton needs none of
these.

\paragraph{The dense shape block.} The shape block is low-dimensional
(interface harmonics up to some degree on a few interfaces), so the
Gauss--Newton matrix is built explicitly: one Poisson and one saddle
solve per parameter give $\partial r[\bm{\theta}_i]$, and
$H_{\mathrm{GN},ij} = \ip{\partial r[\bm{\theta}_i]}{\partial
r[\bm{\theta}_j]}$, the field-level sensitivities paired on both sides
with lifted basis directions. Away from $J = 0$ these entries depend on
the lift at $O(\norm{r}) = O(J^{1/2})$, because $r$ is transported, not
invariant, under relabelling. That is the same order as the neglected
residual term, and the dependence disappears on the zero set. The
gradient is checked coefficient by coefficient against central
differences of $J$, and the sensitivities $\delta\bu_{\bm{\theta}_i}$
against differences of the state, using the interface-shift machinery
of the perturbation benchmark: the same validation pattern, applied to
the optimisation.

\paragraph{Joint steps.} Density and shape derivatives of the residual
are columns of one linear map $(\delta\rho,\delta s)\mapsto\partial r$, so
the joint Gauss--Newton matrix
$\begin{bmatrix}H_{\rho\rho} & H_{\rho s}\\ H_{s\rho} & H_{ss}\end{bmatrix}$
is a Gram matrix and positive semi-definite, and a joint Newton step is
coherent. The mixed product $H_{s\rho}d = (\ip{\partial r[\bm{\theta}_i]}
{\partial r[d]})_i$ costs one density sensitivity. The step is computed by
preconditioned CG on the joint system, with block preconditioner
$\operatorname{diag}\bigl(\Pi_GG^{-1}, (H_{ss} + \lambda_sG_s)^{-1}\bigr)$,
the second block dense and small. Alternatively the shape block is
eliminated through a Schur complement, at one inner density solve per
shape parameter.

\subsection{Discrete exactness and verification}

Everything in this section differentiates the \emph{discrete} $J_h$
exactly, provided the derivatives of the forms are those of the
assembled forms: the same mesh, the same quadrature, and integrands
differentiated analytically through $F$, $F^{-1}$, $J$ at the quadrature
points, which do not move. The DtN block never contributes, because the
map is the identity near $\partial D$. The finite-difference harness that
validates gradients therefore validates Hessian products the same way:
\begin{itemize}
\item symmetry, $d_1^{T}(Hd_2) = d_2^{T}(Hd_1)$, for the full and the
  Gauss--Newton products separately;
\item $(j(\rho + \eta d) - j(\rho - \eta d))/2\eta$ against $Hd$, with
  $O(\eta^2)$ convergence;
\item $H_{\mathrm{R}}[d,d] = O(J^{1/2})$ along a converging sequence;
\item for shape, the dense $H_{\mathrm{GN}}$ entries against central
  differences of the residual, and the gauge identity
  \eqref{eq:gauge-identity}.
\end{itemize}

% ======================================================================
\section{Algorithms, implementation and measured behaviour}
\label{sec:algorithms}

This section describes the stage-1 machinery as implemented
(\texttt{equilibrium\_figures.hpp}, \texttt{riesz.hpp},
\texttt{descent.hpp}, \texttt{examples/equilibrium\_density.cpp}). It
explains each algorithm from its model problem, because the behaviour
measured on the first test problem is only intelligible through that
model.

\subsection{Persistent solvers and the cost unit}

In stage~1 no operator depends on the density (\S\ref{sec:hessian}). The
Poisson system (Laplacian plus DtN on the ball) and the Taylor--Hood
saddle are therefore assembled and preconditioned once
(\texttt{DensityFeasibilityProblem}, with its internal Poisson and
\texttt{StokesSaddleSolver} objects), and every evaluation, gradient and
Hessian product only assembles right-hand sides and calls the same
solvers. The unit of cost is a solve:
\begin{center}
\begin{tabular}{l c c}
  & saddle & Poisson \\ \hline
  value $J$ & 1 & 1 \\
  value and gradient & 1 & 2 \\
  Gauss--Newton product & 1 & 2 \\
  full Hessian product (as implemented) & 1 & 3
\end{tabular}
\end{center}
The Gauss--Newton product has a cheap assembly because of an identity.
By the saddle relation of \S\ref{sec:hessian},
$\delta\bu_2^{T}\mathsf{A}\,\delta\bu_1 = -\delta\bu_2^{T}\partial
\mathsf{f}[d_1]$: the Gauss--Newton dual is the \emph{gradient assembly}
\eqref{eq:gradient} with $(\bu, w)$ replaced by the sensitivities
$(\delta\bu, \delta w)$, $\delta w$ being the Poisson adjoint sourced by
$\rho\,\delta\bu$. The residual term needs one more Poisson solve (the
adjoint sourced by $d\,\bu$). By linearity the two adjoints can be merged
into one solve, as in \eqref{eq:hessvec}, which would bring the full
product to the Gauss--Newton cost.

\subsection{Metric quantities without Gram matrices}

The Riesz maps of \S\ref{sec:sobolev} are implemented as operators taking
a dual $j$ to a gradient $g$: the $L^2$ map (a mass solve), and the
Sobolev map $\tilde A = \alpha M + \beta K$ with a homogeneous Dirichlet
condition on the outer boundary of the extension domain, applied $s$
times with mass products between ($s = 2$ the default). Vector-valued
spaces get the component-wise metric, for later use with mapping
controls. No Gram inner product is ever assembled. Every metric quantity
the loops need is a pairing of a dual with a field, since
$\ip{g}{h}_G = (Gg)^{T}h = j^{T}h$ for $g = G^{-1}j$: the squared gradient
norm is $j^{T}g$, the slope along $d$ is $j^{T}d$, and the mass
projection scalar is $\mathsf{m}^{T}g$.

\subsection{Projected nonlinear conjugate gradients}

\paragraph{The constrained gradient.} Let $T = \{\delta : \mathsf{m}^{T}\delta
= 0\}$ be the fixed-mass plane, $r = G^{-1}\mathsf{m}$ the cached
representative of the mass functional, and
\begin{equation}
  \Pi = I - \frac{r\,\mathsf{m}^{T}}{\mathsf{m}^{T}r} .
\end{equation}
$\Pi$ is idempotent ($\Pi r = 0$, $\Pi\delta = \delta$ on $T$) and
$G$-self-adjoint: $G\Pi = G - \mathsf{m}\mathsf{m}^{T}/\mathsf{m}^{T}r$ is
symmetric, using $Gr = \mathsf{m}$. Hence the projected gradient
$g_\Pi = \Pi G^{-1}j$ is the Riesz representative of $j$ \emph{restricted
to} $T$: for $\delta\in T$, $\ip{g_\Pi}{\delta}_G = j^{T}G^{-1}(G\Pi)\delta
= j^{T}\Pi\delta = j^{T}\delta$. Pairings of raw duals with projected
fields are therefore the inner products of the constrained problem, e.g.\
$\norm{g_\Pi}_G^2 = j^{T}g_\Pi$.

\paragraph{Why the gradient, not the direction, is projected.} At a
constrained minimiser $j = \mu\,\mathsf{m}$ (the mass multiplier), so the
unprojected gradient tends to $\mu r \ne 0$. A loop that projects only
its search direction computes the Polak--Ribi\`ere coefficient
\begin{equation}
  \beta_k = \max\Bigl(0,\ \frac{j_k^{T}(g_k - g_{k-1})}{j_{k-1}^{T}g_{k-1}}\Bigr)
\end{equation}
from gradients dominated, near convergence, by this non-vanishing normal
component. The ratio then measures the wrong thing and the conjugacy is
lost. With $g = g_\Pi$ throughout, $\beta_k$ is the Polak--Ribi\`ere
coefficient of the problem on $T$, and the method is genuinely
conjugate-gradient-like there. In the measured runs this change alone
cut the iteration count by roughly an order of magnitude.
$d_k = -g_k + \beta_kd_{k-1}$ stays in $T$. If $j_k^{T}d_k \ge 0$ (a stale
direction), the loop restarts with $-g_k$; if the restarted direction does
not descend either, the constrained minimum has been reached to solver
accuracy.

\paragraph{Line search.} Along a ray $\rho + t d$ the functional is an
exact quartic in $t$ (\S\ref{sec:derivative}; the quadratic prior keeps it
quartic). The implemented search is \emph{forward tracking}: start from
the previous accepted step (warm start), halve until $J$ decreases, then
double while it keeps decreasing. If the doubling phase advanced at least
once and the quartic is unimodal on the ray, the accepted $t$ satisfies
$J(t/2) > J(t) \le J(2t)$, which brackets the ray minimiser in
$(t/2, 2t)$: within a factor two. If the halving phase ends the search,
only $t^\star < 2t$ is guaranteed. Acceptance is on simple decrease; the
exact quartic line search of \S\ref{sec:second} (one Poisson and two
saddle solves) is the available upgrade.

\subsection{Levenberg--Marquardt Gauss--Newton}
\label{sec:lm}

\paragraph{The Gauss--Newton model.} $J = \tfrac12\norm{r(\rho)}^2$ is a
nonlinear least-squares problem with residual $r = (2\mu)^{1/2}\eps(\bu)$
(\S\ref{sec:hessian}). Linearising the residual, $r(\rho + s) \approx r +
Ls$ with $L = \partial r$, gives the quadratic model
\begin{equation}
  m(s) = \tfrac12\norm{r + Ls}^2
       = J + j^{T}s + \tfrac12 s^{T}H_{\mathrm{GN}}s,
  \qquad j = L^{T}r,\quad H_{\mathrm{GN}} = L^{T}L .
\end{equation}
Minimising the model gives the Gauss--Newton step $H_{\mathrm{GN}}s = -j$.
It uses only first derivatives of the residual, yet near a zero-residual
solution it behaves like Newton's method, because the neglected Hessian
term is $O(\norm{r})$ (\S\ref{sec:hessian}). Taken raw, it is useless
here: $H_{\mathrm{GN}}$ is compact (order $-2$) and exactly singular along
the barotropic family, so $s$ is dominated by the poorly determined
directions and leaves the region where the model is valid.

\paragraph{Damping: two equivalent views.} Levenberg
\cite{Levenberg1944} and Marquardt \cite{Marquardt1963} replace the step
by
\begin{equation}
  (H_{\mathrm{GN}} + \lambda D)\,s = -j ,
  \label{eq:lm-step}
\end{equation}
with $D$ symmetric positive definite (implemented: the mass matrix $M$ of
the control space) and $\lambda > 0$.
\begin{itemize}
\item \emph{Trust region.} \eqref{eq:lm-step} with
  $\lambda \ge 0$, $\lambda(\norm{s}_D - \Delta) = 0$ is exactly the
  optimality system of $\min m(s)$ subject to $\norm{s}_D \le \Delta$
  (Mor\'e \cite{More1978}). $\lambda$ is the multiplier of the trust
  region: large $\lambda$ means a small region. As $\lambda\to\infty$,
  $s \approx -\lambda^{-1}D^{-1}j$ is a short steepest-descent step in the
  $D$-metric; as $\lambda\to0$ it is the Gauss--Newton step. LM
  interpolates between the robust and the fast method.
\item \emph{Tikhonov on the step.} \eqref{eq:lm-step} minimises
  $\tfrac12\norm{r + Ls}^2 + \tfrac\lambda2\norm{s}_D^2$: each step is a
  Tikhonov-regularised solution of the linearised problem.
\end{itemize}

\paragraph{Damping as iterated Tikhonov, and semi-convergence.} For a
linear residual $r(\rho) = r_0 + L\rho$, repeated LM steps with fixed
$\lambda$ are $\rho_{k+1} = \rho_k - (L^TL + \lambda D)^{-1}L^Tr(\rho_k)$.
That is the stationary iterated Tikhonov iteration of
\texttt{doc/gauge\_penalty\_iteration.tex}, with $A = L^TL$ and $Q = D$.
In the generalised singular basis of $(L, D)$, with values $\sigma_i$, the
component of the least-squares error along mode $i$ contracts per step by
\begin{equation}
  \frac{\lambda}{\sigma_i^2 + \lambda} .
\end{equation}
Modes with $\sigma_i^2 \gg \lambda$ are fitted at once; modes with
$\sigma_i^2 \ll \lambda/k$ are still essentially untouched after $k$
steps. The damping is thus the regularisation parameter of an iterative
regularisation method, with effective strength $\lambda/k$ (Hanke
\cite{Hanke1997} analyses LM in exactly this role).

This matters because the discrete functional $J_h$ is not the continuum
$J$. Its small-$\sigma$ modes include directions that lower $J_h$ by
fitting the \emph{discretisation} imbalance rather than the physical one:
densities whose discrete force happens to be nearly a discrete gradient.
These are rough and unphysical. Fitting the physical imbalance first and
the discretisation error later is semi-convergence: the distance to the
physical solution falls and then rises, while $J_h$ falls throughout.
Damping is what holds the late modes back. Letting $\lambda\to0$ removes
the protection, and the iteration then overfits like any unregularised
method. The remedies are those of regularisation theory:
\begin{itemize}
\item stop at the discretisation floor (the discrepancy principle,
  \cite{EnglHankeNeubauer}), estimated from the unperturbed model or by
  refinement, or located without being known by the stagnation of the
  dimensionless infeasibility (\S\ref{sec:stopping});
\item add a prior that changes the problem into a well-posed one, after
  which $\lambda\to0$ is safe and the iteration converges \emph{onto} the
  regularised minimiser.
\end{itemize}
The implemented prior is the roughness of the correction,
$\tfrac{\lambda_p}2(\rho - \rho_0)^{T}K(\rho - \rho_0)$ with $K$ the
$H^1$-seminorm Gram and $\rho_0$ the starting density. It annihilates
constants, so it does not fight the mass constraint, and among
near-equivalent densities it selects the smoothest correction. Its Hessian
$\lambda_pK$ is exact and is added to the model, so the inner system is
\begin{equation}
  (H_{\mathrm{GN}} + \lambda M + \lambda_p K)\,s = -(j + \lambda_pK(\rho - \rho_0)) .
\end{equation}

\paragraph{Choice of prior.} The $H^1$-seminorm prior selects the smooth
member of the near-null family, and is the right default when the
disequilibrium is smooth and so is the correction it requires. It is not
a universal safeguard. When the required correction is oscillatory ---
the density rearrangement that follows a wavy fluid--solid interface
topography, whose equipotentials undulate with the interface --- the
prior penalises exactly the correction that is needed and blocks
convergence entirely. There the Levenberg--Marquardt damping alone must
serve as the semi-convergence protection, with the stopping rule of
\S\ref{sec:stopping} ending the iteration at the floor. The choice of
prior is a statement about the expected correction.

\paragraph{Adapting the damping.} The implemented rule is multiplicative:
solve, project the step onto the mass plane, evaluate; on a decrease of
the objective accept and set $\lambda\leftarrow\lambda/3$, otherwise
reject, set $\lambda\leftarrow4\lambda$ and re-solve (at most eight tries).
This is a poor man's trust region. A proper trust-region method compares
the actual reduction with the reduction predicted by the model,
$\varrho = (J(\rho) - J(\rho+s))/(m(0) - m(s))$, and enlarges or shrinks
the region according to $\varrho$ \cite{More1978,NocedalWright}. The
multiplicative rule uses only the sign of the actual reduction. Its
standard behaviour:
\begin{itemize}
\item $\lambda$ is large far from the solution (short, gradient-like
  steps) and decays geometrically once the model is trustworthy
  (Gauss--Newton steps);
\item the asymmetric factors make a reject--accept cycle increase
  $\lambda$ by $4/3$, so oscillation drifts toward caution rather than
  cycling.
\end{itemize}
The ratio test is the refinement to adopt if step rejections become
frequent.

\paragraph{Inexact inner solves.} \eqref{eq:lm-step} is solved by
preconditioned CG with matrix-free products, to relative tolerance
$10^{-2}$ and at most 30 iterations: a truncated (inexact) Newton method
\cite{DemboEisenstatSteihaug1982}. The loose tolerance costs nothing:
\begin{itemize}
\item The model is only a local approximation of $J$; solving it more
  accurately than it approximates $J$ buys nothing.
\item Inexact Newton theory: with relative residual $\eta_k$ the local
  convergence is linear with rate about $\eta_k$, and superlinear if
  $\eta_k\to0$; $10^{-2}$ is already a fast linear rate.
\item CG from $s = 0$ decreases the model monotonically, and its iterates
  grow monotonically in the energy norm (Steihaug \cite{Steihaug1983}).
  Every truncated iterate is therefore a descent step for the model, and
  early truncation shortens the step --- truncation is itself a
  regularisation, of the same kind as the damping.
\end{itemize}
The step is projected onto the mass plane after the inner solve. This is a
simplification: the exact constrained model step would run the inner CG
with the projected preconditioner of \S\ref{sec:newton-krylov}.

\paragraph{Spectrum and preconditioning.} $H_{\mathrm{GN}} = L^TL$ is
positive semi-definite, so CG is well defined on
\eqref{eq:lm-step} for any $\lambda > 0$. Its convergence is governed by
the spectrum of the preconditioned operator. Implemented, the
preconditioner is the mass solve $M^{-1}$:
\begin{equation}
  M^{-1}(H_{\mathrm{GN}} + \lambda M + \lambda_pK)
  = \lambda I + M^{-1}H_{\mathrm{GN}} + \lambda_p M^{-1}K .
\end{equation}
\begin{itemize}
\item \emph{Without the prior} this is $\lambda I$ plus a compact operator
  (order $-2$). Its eigenvalues cluster at $\lambda$ with finitely many
  outliers above any threshold, independently of the mesh. CG then
  converges superlinearly, with iteration counts governed by the number of
  outliers above about $\lambda$, a number that does not grow under
  refinement but does grow as $\lambda\to0$.
\item \emph{With the prior} the term $\lambda_pM^{-1}K$ is
  \emph{unbounded}. Its eigenvalues reach $O(\lambda_ph^{-2})$, the
  condition number grows like $\lambda_ph^{-2}/\lambda$, and CG
  iterations like $h^{-1}(\lambda_p/\lambda)^{1/2}$: mesh-dependent, and
  worse as $\lambda$ decays. With the iteration cap of 30 this shows up as
  earlier truncation, not divergence. The mesh-independent choice is the
  Riesz map of the regularised metric itself, any fixed Sobolev map
  $(\alpha M + \beta K)^{-1}$ with $\alpha,\beta > 0$. It is spectrally
  equivalent to $\lambda M + \lambda_pK$ with constants
  $\max(\lambda/\alpha, \lambda_p/\beta)/\min(\lambda/\alpha,
  \lambda_p/\beta)$, independent of $h$, and the preconditioned operator
  is that bounded, well-conditioned part plus a compact perturbation from
  $H_{\mathrm{GN}}$.
\end{itemize}

\paragraph{At the floor.} The residual term neglected by Gauss--Newton is
$O(\norm{r}) = O(J^{1/2})$. On the regularised problem the residual at the
minimiser is small but not zero, so Gauss--Newton converges locally
linearly, with a rate proportional to that residual: fast when the
residual is small, quadratic only in the zero-residual limit. With the
prior the model Hessian $H_{\mathrm{GN}} + \lambda_pK$ is non-singular on
the mass plane (the prior controls everything but constants, which the
mass constraint removes). $\lambda$ can then decay to zero safely, and
the iteration converges onto the regularised minimiser in few outer
steps. This is what is measured.

\subsection{The advection loop}

Poisson for $\Phi$, saddle for $\bu$, the weak-form step
\eqref{eq:weak-advection} capped by CFL, acceptance on $E$ with the exact
rate $-2J$, and the absolute ceiling on $J$. Its measured fate and the
structural fix are in \S\ref{sec:advection-caution}.

\subsection{Bounds}

Not implemented. Pointwise bounds $\rho_{\min}\le\rho\le\rho_{\max}$ with
the mass constraint make the projection a singly linearly constrained,
box-constrained quadratic program. MFEM's \texttt{SLBQPOptimizer} solves
$\min\norm{x - x_t}_2$ subject to $w^Tx = a$, $x_{\mathrm{lo}}\le x\le
x_{\mathrm{hi}}$ in the Euclidean coefficient norm, which is the $L^2$
projection for a lumped mass after scaling by $M_L^{1/2}$. Projection with
bounds in a non-diagonal metric is a genuine QP; the practical compromise
is the metric for the direction and lumped-$L^2$ projection for
feasibility. Nodal bounds imply pointwise bounds only at order one (or in
a Bernstein basis).

\subsection{Stopping diagnostic}
\label{sec:stopping}

$J$ is dimensional, and its floor depends on mesh and model. The
implemented diagnostic is the dimensionless global infeasibility
\begin{equation}
  \eta = \frac{\norm{\dev T_{\min}}_{L^2(\Bf)}}{\norm{p}_{L^2(\Bf)}}
       = \frac{(2J)^{1/2}}{\norm{p}_{L^2(\Bf)}}
  \qquad(\mu_{\mathrm{f}} = \tfrac12),
  \label{eq:eta}
\end{equation}
how far from fluid the minimum-deviatoric state is. The loops stop on
it in one of two ways: at an explicit threshold on $\eta$ --- the
discrepancy rule of \S\ref{sec:lm} --- or where $\eta$ stagnates, which
locates the discretisation floor without knowing it. Barotropy itself is
read off the $(\Phi,\rho)$ scatter over the fluid, which collapses onto a
curve at a barotropic state, and its within-equipotential spread.

\paragraph{The pointwise map.} Where the fluid is out of balance is read
from
\begin{equation}
  \eta(\bx) = \frac{|\dev T(\bx)|}{p_{\mathrm{rms}}} ,
  \qquad
  p_{\mathrm{rms}} = \frac{\norm{p - \bar p}_{L^2(\Bf)}}{|\Bf|^{1/2}} ,
\end{equation}
with $\bar p$ the mean of $p$ over $\Bf$: the scale is the RMS pressure
variation, a single gauge-invariant number. The pointwise ratio
$|\dev T|/|p|$ is not usable: an enclosed fluid's pressure is gauged, the
gauged $p$ crosses zero, and the ratio is singular on its zero set,
displaying the gauge rather than the imbalance. In the solver's
mean-zero gauge $\norm{p} = \norm{p - \bar p}$, so the map's $L^2$ norm is
\begin{equation}
  \norm{\eta(\cdot)}_{L^2(\Bf)} = \frac{(2J)^{1/2}}{p_{\mathrm{rms}}}
  = \eta\,|\Bf|^{1/2} ,
\end{equation}
the global $\eta$ of \eqref{eq:eta} up to the volume factor. The solid
components' $|\dev T|$ recovered in \S\ref{sec:recovery} are mapped on
the same scale, so the whole-body picture is comparable across regions:
at a restored state the fluid's map sits at the floor while an
aspherical solid keeps its unavoidable share. Since $\dev T$ in the fluid
is built from $\eps(\bu_h)$, which is discontinuous across elements, the
fluid map is projected elementwise.

\subsection{Verification}
\label{sec:verification}

Implemented tests (\texttt{TestEquilibriumFigures}, serial and parallel):
\begin{itemize}
\item the fluid-only derivative against central differences;
\item the whole-body derivative and the Euler identity \eqref{eq:euler};
\item agreement of the derivative assembled on the parent mesh with the
  one on the Stokes SubMesh;
\item the persistent problem against the one-shot evaluation;
\item Hessian products against differences of the gradient, and their
  symmetry.
\end{itemize}
Still to come:
\begin{itemize}
\item the cross-check of $\Jmu$ against $\JF$ over a contrast sweep
  (Proposition~\ref{prop:Jmu}) and the leakage control $\Ja > 0$
  (Proposition~\ref{prop:leakage});
\item the stage-2 gradient and gauge checks of \S\ref{sec:second}.
\end{itemize}

\subsection{Measured behaviour: density restoration}
\label{sec:measured}

\paragraph{Set-up.} The three-layer model of the test meshes (solid inner
core, fluid outer core, solid mantle, buffer), with fluid density
$\rho_0 = 1.2 - 0.3r^2 + 0.1\,x_2/r$. The base is barotropic; the
degree-one lateral part of amplitude $0.1$ is not. The functional is
$\JF$ with $\bu = 0$ on both fluid boundaries, enriched by the inner
core's force and torque balance through the rigid border of
Remark~\ref{rem:innercore} (the enriched certificate is the
implementation's default; the measurements below that predate the
border were taken with the clamped $\JF$, whose value a restored state
rejoins to within a fraction of a percent).

\paragraph{The geometric family.} Disequilibrium can also be produced by
the geometry alone, without a hand-inserted lateral density term. In an
aspherical container --- flattened interfaces, or topography on the
fluid--solid interface --- the potential is aspherical and its level sets
are not spheres, so a purely radial, non-uniform starting density is not
constant on them and, by Proposition~\ref{prop:barotropic}, the model
has no static state at that density. The signal is linear in the shape
amplitude (the non-barotropic part of the load, the flow $\bu$ and
$J^{1/2}$, hence $\eta$, are all first order in it), so a small amplitude
competes with the discretisation floor, and the discretisation order must
put the floor below the signal.

\paragraph{The floor dominates.} The unperturbed base model, an exact
equilibrium in the continuum, evaluates to $J_h = 9.4\times10^{-9}$. That
is two thirds of the perturbed start, $1.4\times10^{-8}$: most of the
starting value is the discretisation imbalance of the radial state, not
the lateral perturbation. A reduction of $J_h$ by more than a factor of
about $1.5$ is therefore not progress toward the physical answer but
fitting of discretisation error. The success measures are the
within-equipotential density spread, the scatter, and the physical range
of the density; $J_h$ itself is not one.

\paragraph{Results.}
\begin{center}
\begin{tabular}{p{0.3\textwidth} p{0.6\textwidth}}
  run & outcome \\ \hline
  NCG, $L^2$ metric, no prior &
    $J_h$ down by a factor $9$ while the density runs to
    $[-4.5,\ 8.4]$: semi-convergence, deep into the overfitting regime. \\
  NCG, $H^1$ metric &
    iterates stay bounded but roughen. \\
  NCG, $H^2$ metric &
    stays physical; over a 400-iteration run $J_h$ falls below the base
    floor and the within-equipotential spread grows again
    (late-stage semi-convergence). \\
  LM--GN, no prior &
    $J_h$ at $0.236$ of the start in 15 iterations; once $\lambda\to0$ it
    overfits like the unregularised descent. \\
  LM--GN, prior $\lambda_p = 10^{-7}$ &
    the best restoration of the study: spread $3.4\times10^{-2}\to
    1.37\times10^{-2}$, density within $[1.09,\ 1.22]$, in 15 iterations ---
    an order of magnitude fewer than NCG with the same prior --- converging
    onto the regularised solution.
\end{tabular}
\end{center}

\paragraph{Reading.} The Sobolev order is a regularisation that delays
semi-convergence ($L^2$, then $H^1$, then $H^2$, in increasing
robustness) but does not prevent it: a descent on $J_h$ alone eventually
fits the discretisation. The prior changes the problem. With it, the
second-order method does what the theory of \S\ref{sec:lm} predicts:
Gauss--Newton exact at the floor, $\lambda\to0$ safe, and convergence onto
the regularised minimiser in few steps. The remaining spread
($1.37\times10^{-2}$) is the regularised answer at this mesh and prior
weight. Its dependence on $\lambda_p$ and $h$, together with the targets
of the milestone below, is the next study.

\paragraph{Milestone targets.} With a prior centred on the barotropic
base, the optimisation route should recover the base (a zero of $J$ and
of the prior). The implemented roughness prior is centred on the start,
which selects the smoothest feasible correction instead. The advection
route, once range-preserving, should reach a stable rearrangement of the
starting material. Its target is explicit here. The spherical mantle
exerts no force in the cavity (Remark~\ref{rem:shell}). The inner core is
uniform (density $1.3$) and at least as dense as any fluid value. Applying
the Riesz rearrangement inequality to the combined density (inner core
plus fluid), whose symmetric decreasing rearrangement keeps the inner
core in place, gives the global minimiser of $E$ among fluid
rearrangements: the radially decreasing rearrangement $\rho^{*}$ of the
start within the fluid shell. $\rho^{*}$ can be computed directly from the
distribution function. Further studies: iteration counts of the metrics and of LM
under refinement, optimisation against advection, then $\JF$ against
$\Jmu$ on an aspherical model.

% ======================================================================
\begin{thebibliography}{99}

\bibitem{AW10}
D.~Al-Attar, J.~H. Woodhouse,
\emph{On the parametrization of equilibrium stress fields in the Earth},
Geophys. J. Int. 181 (2010), 567--576.

\bibitem{AC16}
D.~Al-Attar, O.~Crawford,
\emph{Particle relabelling transformations in elastodynamics},
Geophys. J. Int. 205 (2016), 575--593.

\bibitem{MA24}
M.~Maitra, D.~Al-Attar,
\emph{On the elastodynamics of rotating planets},
Geophys. J. Int. 237 (2024), 1301--1338.

% TODO: verify page range of Friedman & Schutz (1978).
\bibitem{FriedmanSchutz1978}
J.~L. Friedman, B.~F. Schutz,
\emph{Lagrangian perturbation theory of nonrelativistic fluids},
Astrophys. J. 221 (1978), 937--957.

\bibitem{GR86}
V.~Girault, P.-A. Raviart,
\emph{Finite Element Methods for Navier--Stokes Equations},
Springer, 1986.

\bibitem{TaylorHood}
C.~Taylor, P.~Hood,
\emph{A numerical solution of the Navier--Stokes equations using the
finite element technique},
Comput. Fluids 1 (1973), 73--100.

\bibitem{Danskin}
J.~M. Danskin,
\emph{The Theory of Max-Min and its Application to Weapons Allocation
Problems},
Springer, 1967.

\bibitem{Balakrishnan1960}
A.~V. Balakrishnan,
\emph{Fractional powers of closed operators and the semigroups generated
by them},
Pacific J. Math. 10 (1960), 419--437.

\bibitem{BonitoPasciak2015}
A.~Bonito, J.~E. Pasciak,
\emph{Numerical approximation of fractional powers of elliptic
operators},
Math. Comp. 84 (2015), 2083--2110.

% TODO: verify the precise statement used (maximum of a strictly convex,
% weakly sequentially continuous functional over the weak closure of a
% rearrangement class is attained in the class) and its location in Burton (1987).
\bibitem{Burton1987}
G.~R. Burton,
\emph{Rearrangements of functions, maximization of convex functionals,
and vortex rings},
Math. Ann. 276 (1987), 225--253.

% TODO: verify theorem number of the Riesz rearrangement inequality.
\bibitem{LiebLoss}
E.~H. Lieb, M.~Loss,
\emph{Analysis}, 2nd ed.,
American Mathematical Society, 2001.

% TODO: verify volume/pages.
\bibitem{HiptmairPaganiniSargheini2015}
R.~Hiptmair, A.~Paganini, S.~Sargheini,
\emph{Comparison of approximate shape gradients},
BIT Numer. Math. 55 (2015), 459--485.

% TODO: verify issue/pages.
\bibitem{PolakRibiere1969}
E.~Polak, G.~Ribi\`ere,
\emph{Note sur la convergence de m\'ethodes de directions conjugu\'ees},
Rev. Fran\c{c}aise Informat. Recherche Op\'erationnelle 3 (1969), 35--43.

\bibitem{Steihaug1983}
T.~Steihaug,
\emph{The conjugate gradient method and trust regions in large scale
optimization},
SIAM J. Numer. Anal. 20 (1983), 626--637.

% TODO: verify volume/pages.
\bibitem{GouldHribarNocedal2001}
N.~I.~M. Gould, M.~E. Hribar, J.~Nocedal,
\emph{On the solution of equality constrained quadratic programming
problems arising in optimization},
SIAM J. Sci. Comput. 23 (2001), 1376--1395.

% TODO: verify edition and the chapter on second-order shape derivatives.
\bibitem{DelfourZolesio}
M.~C. Delfour, J.-P. Zol\'esio,
\emph{Shapes and Geometries: Metrics, Analysis, Differential Calculus,
and Optimization}, 2nd ed., SIAM, 2011.

\bibitem{Levenberg1944}
K.~Levenberg,
\emph{A method for the solution of certain non-linear problems in least
squares},
Quart. Appl. Math. 2 (1944), 164--168.

\bibitem{Marquardt1963}
D.~W. Marquardt,
\emph{An algorithm for least-squares estimation of nonlinear parameters},
J. Soc. Indust. Appl. Math. 11 (1963), 431--441.

% TODO: verify pages.
\bibitem{More1978}
J.~J. Mor\'e,
\emph{The Levenberg--Marquardt algorithm: implementation and theory},
in G.~A. Watson (ed.), \emph{Numerical Analysis}, Lecture Notes in
Math. 630, Springer, 1978, 105--116.

% TODO: verify volume/pages.
\bibitem{Hanke1997}
M.~Hanke,
\emph{A regularizing Levenberg--Marquardt scheme, with applications to
inverse groundwater filtration problems},
Inverse Problems 13 (1997), 79--95.

\bibitem{DemboEisenstatSteihaug1982}
R.~S. Dembo, S.~C. Eisenstat, T.~Steihaug,
\emph{Inexact Newton methods},
SIAM J. Numer. Anal. 19 (1982), 400--408.

\bibitem{EnglHankeNeubauer}
H.~W. Engl, M.~Hanke, A.~Neubauer,
\emph{Regularization of Inverse Problems},
Kluwer, 1996.

\bibitem{NocedalWright}
J.~Nocedal, S.~J. Wright,
\emph{Numerical Optimization}, 2nd ed., Springer, 2006.

\bibitem{Armijo1966}
L.~Armijo,
\emph{Minimization of functions having Lipschitz continuous first
partial derivatives},
Pacific J. Math. 16 (1966), 1--3.

\end{thebibliography}

\end{document}
